
\documentclass{article}
\usepackage{colortbl}
\usepackage{makecell}
\usepackage{multirow}
\usepackage{supertabular}

\begin{document}

\newcounter{utterance}

\twocolumn

{ \footnotesize  \setcounter{utterance}{1}
\setlength{\tabcolsep}{0pt}
\begin{supertabular}{c@{$\;$}|p{.15\linewidth}@{}p{.15\linewidth}p{.15\linewidth}p{.15\linewidth}p{.15\linewidth}p{.15\linewidth}}

    \# & $\;$A & \multicolumn{4}{c}{Game Master} & $\;\:$B\\
    \hline 

    \theutterance \stepcounter{utterance}  

    & & \multicolumn{4}{p{0.6\linewidth}}{\cellcolor[rgb]{0.9,0.9,0.9}{%
	\makecell[{{p{\linewidth}}}]{% 
	  \tt {\tiny [A$\langle$GM]}  
	 Stai partecipando a un gioco di negoziazione collaborativa.\\ \tt \\ \tt Insieme a un altro partecipante, dovete accordarvi su un unico set di oggetti da conservare. Ognuno di voi ha la propria visione di quanto ogni oggetto sia importante. Non sai come l'altro partecipante valuta gli oggetti. Inoltre, ti viene fornito lo sforzo che ogni oggetto richiede.\\ \tt Potete accordarvi su un set solo se lo sforzo totale degli oggetti selezionati non supera un limite condiviso:\\ \tt \\ \tt LIMITE = 2307\\ \tt \\ \tt Ecco i valori di sforzo dei singoli oggetti:\\ \tt \\ \tt Sforzo oggetto = {"A21": 147, "B93": 698, "A11": 588, "B33": 43, "C08": 356, "B39": 80, "B96": 95, "C85": 743, "A82": 106, "B31": 308, "B20": 325, "A60": 255, "B09": 276, "C81": 543, "B21": 51}\\ \tt \\ \tt Ecco la tua visione personale sull'importanza di ogni oggetto:\\ \tt \\ \tt Valori di importanza dell'oggetto = {"A21": 138, "B93": 583, "A11": 65, "B33": 262, "C08": 121, "B39": 508, "B96": 461, "C85": 484, "A82": 668, "B31": 389, "B20": 215, "A60": 97, "B09": 500, "C81": 30, "B21": 400}\\ \tt \\ \tt Obiettivo:\\ \tt \\ \tt Il tuo obiettivo è negoziare un set condiviso di oggetti che ti avvantaggi il più possibile (cioè, massimizzi l'importanza totale per TE), rimanendo entro il LIMITEE. Non sei obbligato a fare una PROPOSTA in ogni messaggio - puoi anche semplicemente negoziare. Tutte le tattiche sono permesse!\\ \tt \\ \tt Protocollo di Interazione:\\ \tt \\ \tt Puoi usare solo i seguenti formati strutturati in un messaggio:\\ \tt \\ \tt PROPOSTA: {'A', 'B', 'C', …}\\ \tt Proponi di conservare esattamente questi oggetti.\\ \tt RIFIUTO: {'A', 'B', 'C', …}\\ \tt Rifiuta esplicitamente la proposta dell'avversario.\\ \tt ARGOMENTAZIONE: {'...'}\\ \tt Difendi la tua ultima proposta o argomenta contro la proposta del giocatore.\\ \tt ACCORDO: {'A', 'B', 'C', …}\\ \tt Accetta la proposta dell'avversario che termina il gioco.\\ \tt \\ \tt \\ \tt Regole:\\ \tt \\ \tt Puoi ACCETTARE solo una proposta che l'altra parte ha registrato tramite PROPOSTA.\\ \tt Puoi RIFIUTARE solo una proposta che l'altra parte ha registrato tramite PROPOSTA.\\ \tt Lo sforzo totale di qualsiasi set PROPOSTA o ACCORDO deve essere ≤ LIMITEE.\\ \tt NON rivelare i tuoi punteggi di importanza nascosti.\\ \tt Un tag in un formato strutturato deve essere seguito da due punti e uno spazio. L'argomento deve essere un set python contenente 0 o più stringhe.\\ \tt Quindi, deve essere della forma TAG: {...}\\ \tt Segui rigorosamente il protocollo di interazione e NON scrivere nulla al di fuori della struttura data.\\ \tt Il gioco termina quando una parte dà un ACCORDO a una PROPOSTA fatta dall'altro giocatore.\\ \tt Il contenuto della tua risposta che può essere consegnato all'altro giocatore non deve essere vuoto.\\ \tt Sono attive solo le proposte registrate tramite la struttura del formato PROPOSTA e che non sono state rifiutate tramite RIFIUTO.\\ \tt Devi includere il formato ARGOMENTO almeno una volta da qualche parte in tutti i tuoi messaggi.\\ \tt  
	  } 
	   } 
	   } 
	 & \\ 
 

    \theutterance \stepcounter{utterance}  

    & \multicolumn{4}{p{0.6\linewidth}}{\cellcolor[rgb]{0.8,1,0.9}{%
	\makecell[{{p{\linewidth}}}]{% 
	  \tt {\tiny [A$\rangle$GM]}  
	 PROPOSTA: {'A82', 'B93', 'B39', 'B09', 'C85', 'B21', 'B31', 'B33'}\\ \tt ARGOMENTAZIONE: {"L'insieme proposto include oggetti con elevata importanza rispetto allo sforzo, che assicurano un buon utilizzo del limite condiviso."} 
	  } 
	   } 
	   } 
	 & & \\ 
 

    \theutterance \stepcounter{utterance}  

    & & \multicolumn{4}{p{0.6\linewidth}}{\cellcolor[rgb]{0.9,0.9,0.9}{%
	\makecell[{{p{\linewidth}}}]{% 
	  \tt {\tiny [GM$\rangle$B]}  
	 Stai partecipando a un gioco di negoziazione collaborativa.\\ \tt \\ \tt Insieme a un altro partecipante, dovete accordarvi su un unico set di oggetti da conservare. Ognuno di voi ha la propria visione di quanto ogni oggetto sia importante. Non sai come l'altro partecipante valuta gli oggetti. Inoltre, ti viene fornito lo sforzo che ogni oggetto richiede.\\ \tt Potete accordarvi su un set solo se lo sforzo totale degli oggetti selezionati non supera un limite condiviso:\\ \tt \\ \tt LIMITE = 2307\\ \tt \\ \tt Ecco i valori di sforzo dei singoli oggetti:\\ \tt \\ \tt Sforzo oggetto = {"A21": 147, "B93": 698, "A11": 588, "B33": 43, "C08": 356, "B39": 80, "B96": 95, "C85": 743, "A82": 106, "B31": 308, "B20": 325, "A60": 255, "B09": 276, "C81": 543, "B21": 51}\\ \tt \\ \tt Ecco la tua visione personale sull'importanza di ogni oggetto:\\ \tt \\ \tt Valori di importanza dell'oggetto = {"A82": 30, "B93": 65, "B39": 97, "B09": 121, "C85": 138, "B96": 215, "B21": 262, "B31": 389, "B33": 400, "B20": 461, "A21": 484, "C08": 500, "A60": 508, "A11": 583, "C81": 668}\\ \tt \\ \tt Obiettivo:\\ \tt \\ \tt Il tuo obiettivo è negoziare un set condiviso di oggetti che ti avvantaggi il più possibile (cioè, massimizzi l'importanza totale per TE), rimanendo entro il LIMITEE. Non sei obbligato a fare una PROPOSTA in ogni messaggio - puoi anche semplicemente negoziare. Tutte le tattiche sono permesse!\\ \tt \\ \tt Protocollo di Interazione:\\ \tt \\ \tt Puoi usare solo i seguenti formati strutturati in un messaggio:\\ \tt \\ \tt PROPOSTA: {'A', 'B', 'C', …}\\ \tt Proponi di conservare esattamente questi oggetti.\\ \tt RIFIUTO: {'A', 'B', 'C', …}\\ \tt Rifiuta esplicitamente la proposta dell'avversario.\\ \tt ARGOMENTAZIONE: {'...'}\\ \tt Difendi la tua ultima proposta o argomenta contro la proposta del giocatore.\\ \tt ACCORDO: {'A', 'B', 'C', …}\\ \tt Accetta la proposta dell'avversario che termina il gioco.\\ \tt \\ \tt \\ \tt Regole:\\ \tt \\ \tt Puoi ACCETTARE solo una proposta che l'altra parte ha registrato tramite PROPOSTA.\\ \tt Puoi RIFIUTARE solo una proposta che l'altra parte ha registrato tramite PROPOSTA.\\ \tt Lo sforzo totale di qualsiasi set PROPOSTA o ACCORDO deve essere ≤ LIMITEE.\\ \tt NON rivelare i tuoi punteggi di importanza nascosti.\\ \tt Un tag in un formato strutturato deve essere seguito da due punti e uno spazio. L'argomento deve essere un set python contenente 0 o più stringhe.\\ \tt Quindi, deve essere della forma TAG: {...}\\ \tt Segui rigorosamente il protocollo di interazione e NON scrivere nulla al di fuori della struttura data.\\ \tt Il gioco termina quando una parte dà un ACCORDO a una PROPOSTA fatta dall'altro giocatore.\\ \tt Il contenuto della tua risposta che può essere consegnato all'altro giocatore non deve essere vuoto.\\ \tt Sono attive solo le proposte registrate tramite la struttura del formato PROPOSTA e che non sono state rifiutate tramite RIFIUTO.\\ \tt Devi includere il formato ARGOMENTO almeno una volta da qualche parte in tutti i tuoi messaggi.\\ \tt \\ \tt \\ \tt Ora riceverai il primo messaggio dell'altro giocatore.\\ \tt \\ \tt PROPOSTA: {'A82', 'B93', 'B39', 'B09', 'C85', 'B21', 'B31', 'B33'}\\ \tt ARGOMENTAZIONE: {"L'insieme proposto include oggetti con elevata importanza rispetto allo sforzo, che assicurano un buon utilizzo del limite condiviso."} 
	  } 
	   } 
	   } 
	 & \\ 
 

    \theutterance \stepcounter{utterance}  

    & & & \multicolumn{4}{p{0.6\linewidth}}{\cellcolor[rgb]{1,0.85,0.72}{%
	\makecell[{{p{\linewidth}}}]{% 
	  \tt {\tiny [GM$\langle$B]}  
	 RIFIUTO: {'A82', 'B93', 'B39', 'B09', 'C85', 'B21', 'B31', 'B33'}\\ \tt PROPOSTA: {'C81', 'A11', 'A60', 'C08', 'A21', 'B33', 'B31', 'B21'}\\ \tt ARGOMENTAZIONE: {"La proposta seleziona oggetti con un equilibrio ottimale tra importanza e sforzo, massimizzando il valore totale rispetto al limite condiviso."} 
	  } 
	   } 
	   } 
	  \\ 
 

    \theutterance \stepcounter{utterance}  

    & & \multicolumn{4}{p{0.6\linewidth}}{\cellcolor[rgb]{0.9,0.9,0.9}{%
	\makecell[{{p{\linewidth}}}]{% 
	  \tt {\tiny [A$\langle$GM]}  
	 RIFIUTO: {'A82', 'B93', 'B39', 'B09', 'C85', 'B21', 'B31', 'B33'}\\ \tt PROPOSTA: {'C81', 'A11', 'A60', 'C08', 'A21', 'B33', 'B31', 'B21'}\\ \tt ARGOMENTAZIONE: {"La proposta seleziona oggetti con un equilibrio ottimale tra importanza e sforzo, massimizzando il valore totale rispetto al limite condiviso."} 
	  } 
	   } 
	   } 
	 & \\ 
 

    \theutterance \stepcounter{utterance}  

    & \multicolumn{4}{p{0.6\linewidth}}{\cellcolor[rgb]{0.8,1,0.9}{%
	\makecell[{{p{\linewidth}}}]{% 
	  \tt {\tiny [A$\rangle$GM]}  
	 RIFIUTO: {'C81', 'A11', 'A60', 'C08', 'A21', 'B33', 'B31', 'B21'}\\ \tt PROPOSTA: {'A82', 'B93', 'B39', 'B09', 'B31', 'B33', 'B21', 'C08', 'B20'}\\ \tt ARGOMENTAZIONE: {"Questa proposta include oggetti selezionati per il loro equilibrio tra valore e sforzo, garantendo un utilizzo efficiente del limite condiviso."} 
	  } 
	   } 
	   } 
	 & & \\ 
 

    \theutterance \stepcounter{utterance}  

    & & \multicolumn{4}{p{0.6\linewidth}}{\cellcolor[rgb]{0.9,0.9,0.9}{%
	\makecell[{{p{\linewidth}}}]{% 
	  \tt {\tiny [GM$\rangle$B]}  
	 RIFIUTO: {'C81', 'A11', 'A60', 'C08', 'A21', 'B33', 'B31', 'B21'}\\ \tt PROPOSTA: {'A82', 'B93', 'B39', 'B09', 'B31', 'B33', 'B21', 'C08', 'B20'}\\ \tt ARGOMENTAZIONE: {"Questa proposta include oggetti selezionati per il loro equilibrio tra valore e sforzo, garantendo un utilizzo efficiente del limite condiviso."} 
	  } 
	   } 
	   } 
	 & \\ 
 

    \theutterance \stepcounter{utterance}  

    & & & \multicolumn{4}{p{0.6\linewidth}}{\cellcolor[rgb]{1,0.85,0.72}{%
	\makecell[{{p{\linewidth}}}]{% 
	  \tt {\tiny [GM$\langle$B]}  
	 RIFIUTO: {'A82', 'B93', 'B39', 'B09', 'B31', 'B33', 'B21', 'C08', 'B20'}\\ \tt PROPOSTA: {'A11', 'A60', 'A21', 'C08', 'B31', 'B33', 'B20', 'B96', 'B39'}\\ \tt ARGOMENTAZIONE: {"Questa proposta seleziona oggetti con un equilibrio eccezionale tra valore e sforzo, escludendo elementi di basso potenziale per massimizzare il valore complessivo consentito dal limite condiviso."} 
	  } 
	   } 
	   } 
	  \\ 
 

    \theutterance \stepcounter{utterance}  

    & & \multicolumn{4}{p{0.6\linewidth}}{\cellcolor[rgb]{0.9,0.9,0.9}{%
	\makecell[{{p{\linewidth}}}]{% 
	  \tt {\tiny [A$\langle$GM]}  
	 RIFIUTO: {'A82', 'B93', 'B39', 'B09', 'B31', 'B33', 'B21', 'C08', 'B20'}\\ \tt PROPOSTA: {'A11', 'A60', 'A21', 'C08', 'B31', 'B33', 'B20', 'B96', 'B39'}\\ \tt ARGOMENTAZIONE: {"Questa proposta seleziona oggetti con un equilibrio eccezionale tra valore e sforzo, escludendo elementi di basso potenziale per massimizzare il valore complessivo consentito dal limite condiviso."} 
	  } 
	   } 
	   } 
	 & \\ 
 

    \theutterance \stepcounter{utterance}  

    & \multicolumn{4}{p{0.6\linewidth}}{\cellcolor[rgb]{0.8,1,0.9}{%
	\makecell[{{p{\linewidth}}}]{% 
	  \tt {\tiny [A$\rangle$GM]}  
	 RIFIUTO: {'A11', 'A60', 'A21', 'C08', 'B31', 'B33', 'B20', 'B96', 'B39'}\\ \tt PROPOSTA: {'A82', 'B93', 'B09', 'B39', 'C85', 'B96', 'B21', 'B33', 'A21'}\\ \tt ARGOMENTAZIONE: {"Questo set include gli oggetti con i valori di importanza più elevati che rimangono entro il limite di sforzo, massimizzando il mio valore totale."} 
	  } 
	   } 
	   } 
	 & & \\ 
 

    \theutterance \stepcounter{utterance}  

    & & \multicolumn{4}{p{0.6\linewidth}}{\cellcolor[rgb]{0.9,0.9,0.9}{%
	\makecell[{{p{\linewidth}}}]{% 
	  \tt {\tiny [GM$\rangle$B]}  
	 RIFIUTO: {'A11', 'A60', 'A21', 'C08', 'B31', 'B33', 'B20', 'B96', 'B39'}\\ \tt PROPOSTA: {'A82', 'B93', 'B09', 'B39', 'C85', 'B96', 'B21', 'B33', 'A21'}\\ \tt ARGOMENTAZIONE: {"Questo set include gli oggetti con i valori di importanza più elevati che rimangono entro il limite di sforzo, massimizzando il mio valore totale."} 
	  } 
	   } 
	   } 
	 & \\ 
 

    \theutterance \stepcounter{utterance}  

    & & & \multicolumn{4}{p{0.6\linewidth}}{\cellcolor[rgb]{1,0.85,0.72}{%
	\makecell[{{p{\linewidth}}}]{% 
	  \tt {\tiny [GM$\langle$B]}  
	 RIFIUTO: {'A82', 'B93', 'B09', 'B39', 'C85', 'B96', 'B21', 'B33', 'A21'}\\ \tt PROPOSTA: {'C81', 'A11', 'A60', 'C08', 'B21', 'B96', 'B33', 'B39', 'B09'}\\ \tt ARGOMENTAZIONE: {"L'insieme selezionato include oggetti con il miglior equilibrio tra importanza e sforzo, evitando quelli che consumano troppe risorse per un valore insufficiente."} 
	  } 
	   } 
	   } 
	  \\ 
 

    \theutterance \stepcounter{utterance}  

    & & \multicolumn{4}{p{0.6\linewidth}}{\cellcolor[rgb]{0.9,0.9,0.9}{%
	\makecell[{{p{\linewidth}}}]{% 
	  \tt {\tiny [A$\langle$GM]}  
	 RIFIUTO: {'A82', 'B93', 'B09', 'B39', 'C85', 'B96', 'B21', 'B33', 'A21'}\\ \tt PROPOSTA: {'C81', 'A11', 'A60', 'C08', 'B21', 'B96', 'B33', 'B39', 'B09'}\\ \tt ARGOMENTAZIONE: {"L'insieme selezionato include oggetti con il miglior equilibrio tra importanza e sforzo, evitando quelli che consumano troppe risorse per un valore insufficiente."} 
	  } 
	   } 
	   } 
	 & \\ 
 

    \theutterance \stepcounter{utterance}  

    & \multicolumn{4}{p{0.6\linewidth}}{\cellcolor[rgb]{0.8,1,0.9}{%
	\makecell[{{p{\linewidth}}}]{% 
	  \tt {\tiny [A$\rangle$GM]}  
	 RIFIUTO: {'C81', 'A11', 'A60', 'C08', 'B21', 'B96', 'B33', 'B39', 'B09'}\\ \tt PROPOSTA: {'A82', 'B93', 'C85', 'B39', 'B09', 'B96', 'B21', 'B33', 'A21'}\\ \tt ARGOMENTAZIONE: {'Questa selezione include gli oggetti con i più alti valori di importanza per me all'interno del limite di sforzo condiviso, massimizzando il mio beneficio totale.'} 
	  } 
	   } 
	   } 
	 & & \\ 
 

    \theutterance \stepcounter{utterance}  

    & & \multicolumn{4}{p{0.6\linewidth}}{\cellcolor[rgb]{0.9,0.9,0.9}{%
	\makecell[{{p{\linewidth}}}]{% 
	  \tt {\tiny [GM$\rangle$B]}  
	 RIFIUTO: {'C81', 'A11', 'A60', 'C08', 'B21', 'B96', 'B33', 'B39', 'B09'}\\ \tt PROPOSTA: {'A82', 'B93', 'C85', 'B39', 'B09', 'B96', 'B21', 'B33', 'A21'}\\ \tt ARGOMENTAZIONE: {'Questa selezione include gli oggetti con i più alti valori di importanza per me all'interno del limite di sforzo condiviso, massimizzando il mio beneficio totale.'} 
	  } 
	   } 
	   } 
	 & \\ 
 

    \theutterance \stepcounter{utterance}  

    & & & \multicolumn{4}{p{0.6\linewidth}}{\cellcolor[rgb]{1,0.85,0.72}{%
	\makecell[{{p{\linewidth}}}]{% 
	  \tt {\tiny [GM$\langle$B]}  
	 RIFIUTO: {'A82', 'B93', 'C85', 'B39', 'B09', 'B96', 'B21', 'B33', 'A21'}\\ \tt PROPOSTA: {'C81', 'A21', 'A60', 'B20', 'C08', 'B33', 'B21', 'B96', 'B31', 'B39'}\\ \tt ARGOMENTAZIONE: {"Questa selezione massimizza il valore rispetto allo sforzo, utilizzando efficientemente il limite condiviso per ottenere il miglior risultato possibile."} 
	  } 
	   } 
	   } 
	  \\ 
 

    \theutterance \stepcounter{utterance}  

    & & \multicolumn{4}{p{0.6\linewidth}}{\cellcolor[rgb]{0.9,0.9,0.9}{%
	\makecell[{{p{\linewidth}}}]{% 
	  \tt {\tiny [A$\langle$GM]}  
	 RIFIUTO: {'A82', 'B93', 'C85', 'B39', 'B09', 'B96', 'B21', 'B33', 'A21'}\\ \tt PROPOSTA: {'C81', 'A21', 'A60', 'B20', 'C08', 'B33', 'B21', 'B96', 'B31', 'B39'}\\ \tt ARGOMENTAZIONE: {"Questa selezione massimizza il valore rispetto allo sforzo, utilizzando efficientemente il limite condiviso per ottenere il miglior risultato possibile."} 
	  } 
	   } 
	   } 
	 & \\ 
 

    \theutterance \stepcounter{utterance}  

    & \multicolumn{4}{p{0.6\linewidth}}{\cellcolor[rgb]{0.8,1,0.9}{%
	\makecell[{{p{\linewidth}}}]{% 
	  \tt {\tiny [A$\rangle$GM]}  
	 RIFIUTO: {'C81', 'A21', 'A60', 'B20', 'C08', 'B33', 'B21', 'B96', 'B31', 'B39'}\\ \tt ARGOMENTAZIONE: {"La proposta include C81 con importanza 30 per uno sforzo di 543, un rapporto molto sfavorevole. Oggetti come A82 (668 importanza, 106 sforzo) e B39 (508 importanza, 80 sforzo) offrono un valore superiore rispetto allo sforzo e meritano di essere prioritizzati."} 
	  } 
	   } 
	   } 
	 & & \\ 
 

    \theutterance \stepcounter{utterance}  

    & & \multicolumn{4}{p{0.6\linewidth}}{\cellcolor[rgb]{0.9,0.9,0.9}{%
	\makecell[{{p{\linewidth}}}]{% 
	  \tt {\tiny [GM$\rangle$B]}  
	 RIFIUTO: {'C81', 'A21', 'A60', 'B20', 'C08', 'B33', 'B21', 'B96', 'B31', 'B39'}\\ \tt ARGOMENTAZIONE: {"La proposta include C81 con importanza 30 per uno sforzo di 543, un rapporto molto sfavorevole. Oggetti come A82 (668 importanza, 106 sforzo) e B39 (508 importanza, 80 sforzo) offrono un valore superiore rispetto allo sforzo e meritano di essere prioritizzati."} 
	  } 
	   } 
	   } 
	 & \\ 
 

    \theutterance \stepcounter{utterance}  

    & & & \multicolumn{4}{p{0.6\linewidth}}{\cellcolor[rgb]{1,0.85,0.72}{%
	\makecell[{{p{\linewidth}}}]{% 
	  \tt {\tiny [GM$\langle$B]}  
	 PROPOSTA: {'A82', 'B33', 'B21', 'B96', 'C08', 'A60', 'A11', 'B20', 'B31', 'B39'}\\ \tt ARGOMENTAZIONE: {"Questa proposta esclude C81 come suggerito, includendo A82 e altri oggetti strategici che massimizzano il valore per entrambi nel rispetto del limite condiviso. Tutti gli elementi selezionati hanno un ottimo rapporto tra sforzo e valore, evitando di spendere risorse inutili."} 
	  } 
	   } 
	   } 
	  \\ 
 

    \theutterance \stepcounter{utterance}  

    & & \multicolumn{4}{p{0.6\linewidth}}{\cellcolor[rgb]{0.9,0.9,0.9}{%
	\makecell[{{p{\linewidth}}}]{% 
	  \tt {\tiny [A$\langle$GM]}  
	 PROPOSTA: {'A82', 'B33', 'B21', 'B96', 'C08', 'A60', 'A11', 'B20', 'B31', 'B39'}\\ \tt ARGOMENTAZIONE: {"Questa proposta esclude C81 come suggerito, includendo A82 e altri oggetti strategici che massimizzano il valore per entrambi nel rispetto del limite condiviso. Tutti gli elementi selezionati hanno un ottimo rapporto tra sforzo e valore, evitando di spendere risorse inutili."} 
	  } 
	   } 
	   } 
	 & \\ 
 

    \theutterance \stepcounter{utterance}  

    & \multicolumn{4}{p{0.6\linewidth}}{\cellcolor[rgb]{0.8,1,0.9}{%
	\makecell[{{p{\linewidth}}}]{% 
	  \tt {\tiny [A$\rangle$GM]}  
	 PROPOSTA: {'B21', 'B39', 'A82', 'B33', 'B96', 'B09', 'A21', 'B31', 'B93', 'B20'}\\ \tt ARGOMENTAZIONE: {"Questa proposta massimizza il valore relativo allo sforzo, scegliendo oggetti con ottimo ritorno come B21, B39 e A82, garantendo un utilizzo efficiente del limite condiviso."} 
	  } 
	   } 
	   } 
	 & & \\ 
 

    \theutterance \stepcounter{utterance}  

    & & \multicolumn{4}{p{0.6\linewidth}}{\cellcolor[rgb]{0.9,0.9,0.9}{%
	\makecell[{{p{\linewidth}}}]{% 
	  \tt {\tiny [GM$\rangle$B]}  
	 PROPOSTA: {'B21', 'B39', 'A82', 'B33', 'B96', 'B09', 'A21', 'B31', 'B93', 'B20'}\\ \tt ARGOMENTAZIONE: {"Questa proposta massimizza il valore relativo allo sforzo, scegliendo oggetti con ottimo ritorno come B21, B39 e A82, garantendo un utilizzo efficiente del limite condiviso."} 
	  } 
	   } 
	   } 
	 & \\ 
 

    \theutterance \stepcounter{utterance}  

    & & & \multicolumn{4}{p{0.6\linewidth}}{\cellcolor[rgb]{1,0.85,0.72}{%
	\makecell[{{p{\linewidth}}}]{% 
	  \tt {\tiny [GM$\langle$B]}  
	 RIFIUTO: {'B21', 'B39', 'A82', 'B33', 'B96', 'B09', 'A21', 'B31', 'B93', 'B20'}\\ \tt PROPOSTA: {'B33', 'B21', 'B96', 'A21', 'A60', 'B20', 'C08', 'B31', 'B39', 'A11'}\\ \tt ARGOMENTAZIONE: {"La selezione di oggetti è ottimizzata per un rapporto importanza/sforzo bilanciato, evitando quelli come A82 e B93 che non contribuiscono in modo efficiente al valore totale consentito dal limite."} 
	  } 
	   } 
	   } 
	  \\ 
 

    \theutterance \stepcounter{utterance}  

    & & \multicolumn{4}{p{0.6\linewidth}}{\cellcolor[rgb]{0.9,0.9,0.9}{%
	\makecell[{{p{\linewidth}}}]{% 
	  \tt {\tiny [A$\langle$GM]}  
	 RIFIUTO: {'B21', 'B39', 'A82', 'B33', 'B96', 'B09', 'A21', 'B31', 'B93', 'B20'}\\ \tt PROPOSTA: {'B33', 'B21', 'B96', 'A21', 'A60', 'B20', 'C08', 'B31', 'B39', 'A11'}\\ \tt ARGOMENTAZIONE: {"La selezione di oggetti è ottimizzata per un rapporto importanza/sforzo bilanciato, evitando quelli come A82 e B93 che non contribuiscono in modo efficiente al valore totale consentito dal limite."} 
	  } 
	   } 
	   } 
	 & \\ 
 

    \theutterance \stepcounter{utterance}  

    & \multicolumn{4}{p{0.6\linewidth}}{\cellcolor[rgb]{0.8,1,0.9}{%
	\makecell[{{p{\linewidth}}}]{% 
	  \tt {\tiny [A$\rangle$GM]}  
	 PROPOSTA: {'A82', 'B39', 'B09', 'B21', 'B33', 'B96', 'B31', 'B20', 'C08', 'A21', 'A60'}\\ \tt ARGOMENTAZIONE: {"Questa proposta include A82 per il massimo valore individuale, insieme a C08 e A60 che l'altro giocatore ha già proposto, bilanciando sforzo e importanza per entrambe le parti."} 
	  } 
	   } 
	   } 
	 & & \\ 
 

    \theutterance \stepcounter{utterance}  

    & & \multicolumn{4}{p{0.6\linewidth}}{\cellcolor[rgb]{0.9,0.9,0.9}{%
	\makecell[{{p{\linewidth}}}]{% 
	  \tt {\tiny [GM$\rangle$B]}  
	 PROPOSTA: {'A82', 'B39', 'B09', 'B21', 'B33', 'B96', 'B31', 'B20', 'C08', 'A21', 'A60'}\\ \tt ARGOMENTAZIONE: {"Questa proposta include A82 per il massimo valore individuale, insieme a C08 e A60 che l'altro giocatore ha già proposto, bilanciando sforzo e importanza per entrambe le parti."} 
	  } 
	   } 
	   } 
	 & \\ 
 

    \theutterance \stepcounter{utterance}  

    & & & \multicolumn{4}{p{0.6\linewidth}}{\cellcolor[rgb]{1,0.85,0.72}{%
	\makecell[{{p{\linewidth}}}]{% 
	  \tt {\tiny [GM$\langle$B]}  
	 RIFIUTO: {'A82', 'B39', 'B09', 'B21', 'B33', 'B96', 'B31', 'B20', 'C08', 'A21', 'A60'}\\ \tt PROPOSTA: {'A11', 'B21', 'B33', 'B96', 'A21', 'A60', 'B20', 'C08', 'B31', 'B39'}\\ \tt ARGOMENTAZIONE: {"La selezione ottimizza il rapporto tra sforzo e valore escludendo oggetti poco efficienti come A82 e includendo A11 per massimizzare il beneficio totale."} 
	  } 
	   } 
	   } 
	  \\ 
 

    \theutterance \stepcounter{utterance}  

    & & \multicolumn{4}{p{0.6\linewidth}}{\cellcolor[rgb]{0.9,0.9,0.9}{%
	\makecell[{{p{\linewidth}}}]{% 
	  \tt {\tiny [A$\langle$GM]}  
	 RIFIUTO: {'A82', 'B39', 'B09', 'B21', 'B33', 'B96', 'B31', 'B20', 'C08', 'A21', 'A60'}\\ \tt PROPOSTA: {'A11', 'B21', 'B33', 'B96', 'A21', 'A60', 'B20', 'C08', 'B31', 'B39'}\\ \tt ARGOMENTAZIONE: {"La selezione ottimizza il rapporto tra sforzo e valore escludendo oggetti poco efficienti come A82 e includendo A11 per massimizzare il beneficio totale."} 
	  } 
	   } 
	   } 
	 & \\ 
 

    \theutterance \stepcounter{utterance}  

    & \multicolumn{4}{p{0.6\linewidth}}{\cellcolor[rgb]{0.8,1,0.9}{%
	\makecell[{{p{\linewidth}}}]{% 
	  \tt {\tiny [A$\rangle$GM]}  
	 PROPOSTA: {'A82', 'B39', 'B21', 'B33', 'B96', 'B09', 'B93', 'B31', 'A21', 'C08'}\\ \tt ARGOMENTAZIONE: {"Questa proposta massimizza il valore totale per me includendo oggetti altamente efficienti come A82 (668 importanza, 106 sforzo) e B93 (583 importanza), ottimizzando il rapporto sforzo/importanza rispetto al limite condiviso."} 
	  } 
	   } 
	   } 
	 & & \\ 
 

    \theutterance \stepcounter{utterance}  

    & & \multicolumn{4}{p{0.6\linewidth}}{\cellcolor[rgb]{0.9,0.9,0.9}{%
	\makecell[{{p{\linewidth}}}]{% 
	  \tt {\tiny [GM$\rangle$B]}  
	 PROPOSTA: {'A82', 'B39', 'B21', 'B33', 'B96', 'B09', 'B93', 'B31', 'A21', 'C08'}\\ \tt ARGOMENTAZIONE: {"Questa proposta massimizza il valore totale per me includendo oggetti altamente efficienti come A82 (668 importanza, 106 sforzo) e B93 (583 importanza), ottimizzando il rapporto sforzo/importanza rispetto al limite condiviso."} 
	  } 
	   } 
	   } 
	 & \\ 
 

    \theutterance \stepcounter{utterance}  

    & & & \multicolumn{4}{p{0.6\linewidth}}{\cellcolor[rgb]{1,0.85,0.72}{%
	\makecell[{{p{\linewidth}}}]{% 
	  \tt {\tiny [GM$\langle$B]}  
	 RIFIUTO: {'A82', 'B39', 'B21', 'B33', 'B96', 'B09', 'B93', 'B31', 'A21', 'C08'}\\ \tt PROPOSTA: {'A11', 'A60', 'C08', 'A21', 'B20', 'B33', 'B31', 'B21', 'B39', 'B96'}\\ \tt ARGOMENTAZIONE: {"Questa selezione ottimizza l'allocazione dello sforzo, privilegiando oggetti con un'elevata redditività rispetto al consumo delle risorse disponibili."} 
	  } 
	   } 
	   } 
	  \\ 
 

    \theutterance \stepcounter{utterance}  

    & & \multicolumn{4}{p{0.6\linewidth}}{\cellcolor[rgb]{0.9,0.9,0.9}{%
	\makecell[{{p{\linewidth}}}]{% 
	  \tt {\tiny [A$\langle$GM]}  
	 RIFIUTO: {'A82', 'B39', 'B21', 'B33', 'B96', 'B09', 'B93', 'B31', 'A21', 'C08'}\\ \tt PROPOSTA: {'A11', 'A60', 'C08', 'A21', 'B20', 'B33', 'B31', 'B21', 'B39', 'B96'}\\ \tt ARGOMENTAZIONE: {"Questa selezione ottimizza l'allocazione dello sforzo, privilegiando oggetti con un'elevata redditività rispetto al consumo delle risorse disponibili."} 
	  } 
	   } 
	   } 
	 & \\ 
 

    \theutterance \stepcounter{utterance}  

    & \multicolumn{4}{p{0.6\linewidth}}{\cellcolor[rgb]{0.8,1,0.9}{%
	\makecell[{{p{\linewidth}}}]{% 
	  \tt {\tiny [A$\rangle$GM]}  
	 RIFIUTO: {'A11', 'A60', 'C08', 'A21', 'B20', 'B33', 'B31', 'B21', 'B39', 'B96'}\\ \tt PROPOSTA: {'A82', 'B39', 'B21', 'B33', 'B96', 'B09', 'B31', 'C08', 'A60', 'B93'}\\ \tt ARGOMENTAZIONE: {"Questa proposta seleziona oggetti con un elevato rapporto tra valore e sforzo, prioritizzando elementi come A82 e B93 che offrono il migliore ritorno per l'utilizzo delle risorse condivise. Il set è ottimizzato per rispettare il limite totale mantenendo un equilibrio efficiente tra tutti gli oggetti considerati."} 
	  } 
	   } 
	   } 
	 & & \\ 
 

    \theutterance \stepcounter{utterance}  

    & & \multicolumn{4}{p{0.6\linewidth}}{\cellcolor[rgb]{0.9,0.9,0.9}{%
	\makecell[{{p{\linewidth}}}]{% 
	  \tt {\tiny [GM$\rangle$B]}  
	 RIFIUTO: {'A11', 'A60', 'C08', 'A21', 'B20', 'B33', 'B31', 'B21', 'B39', 'B96'}\\ \tt PROPOSTA: {'A82', 'B39', 'B21', 'B33', 'B96', 'B09', 'B31', 'C08', 'A60', 'B93'}\\ \tt ARGOMENTAZIONE: {"Questa proposta seleziona oggetti con un elevato rapporto tra valore e sforzo, prioritizzando elementi come A82 e B93 che offrono il migliore ritorno per l'utilizzo delle risorse condivise. Il set è ottimizzato per rispettare il limite totale mantenendo un equilibrio efficiente tra tutti gli oggetti considerati."} 
	  } 
	   } 
	   } 
	 & \\ 
 

    \theutterance \stepcounter{utterance}  

    & & & \multicolumn{4}{p{0.6\linewidth}}{\cellcolor[rgb]{1,0.85,0.72}{%
	\makecell[{{p{\linewidth}}}]{% 
	  \tt {\tiny [GM$\langle$B]}  
	 RIFIUTO: {'A82', 'B39', 'B21', 'B33', 'B96', 'B09', 'B31', 'C08', 'A60', 'B93'}\\ \tt PROPOSTA: {'B33', 'B21', 'A21', 'B96', 'A60', 'B20', 'C08', 'B31', 'B39', 'A11'}\\ \tt ARGOMENTAZIONE: {"Questa selezione esclude oggetti con sforzo elevato rispetto al valore, come A82 e B93, e privilegia quelli con il miglior equilibrio tra costo e ritorno. Il set proposto ottimizza l'utilizzo del limite per massimizzare il valore complessivo."} 
	  } 
	   } 
	   } 
	  \\ 
 

    \theutterance \stepcounter{utterance}  

    & & \multicolumn{4}{p{0.6\linewidth}}{\cellcolor[rgb]{0.9,0.9,0.9}{%
	\makecell[{{p{\linewidth}}}]{% 
	  \tt {\tiny [A$\langle$GM]}  
	 RIFIUTO: {'A82', 'B39', 'B21', 'B33', 'B96', 'B09', 'B31', 'C08', 'A60', 'B93'}\\ \tt PROPOSTA: {'B33', 'B21', 'A21', 'B96', 'A60', 'B20', 'C08', 'B31', 'B39', 'A11'}\\ \tt ARGOMENTAZIONE: {"Questa selezione esclude oggetti con sforzo elevato rispetto al valore, come A82 e B93, e privilegia quelli con il miglior equilibrio tra costo e ritorno. Il set proposto ottimizza l'utilizzo del limite per massimizzare il valore complessivo."} 
	  } 
	   } 
	   } 
	 & \\ 
 

    \theutterance \stepcounter{utterance}  

    & \multicolumn{4}{p{0.6\linewidth}}{\cellcolor[rgb]{0.8,1,0.9}{%
	\makecell[{{p{\linewidth}}}]{% 
	  \tt {\tiny [A$\rangle$GM]}  
	 RIFIUTO: {'A11', 'B21', 'B33', 'B96', 'A21', 'A60', 'B20', 'C08', 'B31', 'B39'}\\ \tt PROPOSTA: {'A82', 'B39', 'B21', 'B96', 'B33', 'B31', 'C08', 'B20', 'B09'}\\ \tt ARGOMENTAZIONE: {"Questa proposta include oggetti come A82 (massimo valore per costi ridotti), B39 e B09 con ottimo rapporto importanza/sforzo, ottimizzando il valore totale per entrambe le parti senza superare il limite."} 
	  } 
	   } 
	   } 
	 & & \\ 
 

    \theutterance \stepcounter{utterance}  

    & & \multicolumn{4}{p{0.6\linewidth}}{\cellcolor[rgb]{0.9,0.9,0.9}{%
	\makecell[{{p{\linewidth}}}]{% 
	  \tt {\tiny [GM$\rangle$B]}  
	 RIFIUTO: {'A11', 'B21', 'B33', 'B96', 'A21', 'A60', 'B20', 'C08', 'B31', 'B39'}\\ \tt PROPOSTA: {'A82', 'B39', 'B21', 'B96', 'B33', 'B31', 'C08', 'B20', 'B09'}\\ \tt ARGOMENTAZIONE: {"Questa proposta include oggetti come A82 (massimo valore per costi ridotti), B39 e B09 con ottimo rapporto importanza/sforzo, ottimizzando il valore totale per entrambe le parti senza superare il limite."} 
	  } 
	   } 
	   } 
	 & \\ 
 

    \theutterance \stepcounter{utterance}  

    & & & \multicolumn{4}{p{0.6\linewidth}}{\cellcolor[rgb]{1,0.85,0.72}{%
	\makecell[{{p{\linewidth}}}]{% 
	  \tt {\tiny [GM$\langle$B]}  
	 PROPOSTA: {'C81', 'A11', 'A60', 'C08', 'A21', 'B96', 'B39', 'B33'}\\ \tt ARGOMENTAZIONE: {"La selezione ottimale per il limite condiviso, maximizzando il valore totale evitando oggetti inefficaci e prioritizzando quelli con il miglior rapporto tra sforzo e valore."} 
	  } 
	   } 
	   } 
	  \\ 
 

    \theutterance \stepcounter{utterance}  

    & & \multicolumn{4}{p{0.6\linewidth}}{\cellcolor[rgb]{0.9,0.9,0.9}{%
	\makecell[{{p{\linewidth}}}]{% 
	  \tt {\tiny [A$\langle$GM]}  
	 PROPOSTA: {'C81', 'A11', 'A60', 'C08', 'A21', 'B96', 'B39', 'B33'}\\ \tt ARGOMENTAZIONE: {"La selezione ottimale per il limite condiviso, maximizzando il valore totale evitando oggetti inefficaci e prioritizzando quelli con il miglior rapporto tra sforzo e valore."} 
	  } 
	   } 
	   } 
	 & \\ 
 

    \theutterance \stepcounter{utterance}  

    & \multicolumn{4}{p{0.6\linewidth}}{\cellcolor[rgb]{0.8,1,0.9}{%
	\makecell[{{p{\linewidth}}}]{% 
	  \tt {\tiny [A$\rangle$GM]}  
	 RIFIUTO: {'C81', 'A11', 'A60', 'C08', 'A21', 'B96', 'B39', 'B33'}\\ \tt ARGOMENTAZIONE: {"La proposta include oggetti con rapporti sforzo/importanza scadenti per me, come C81 (30 su 543) e A11 (65 su 588). Preferisco oggetti con alto rapporto come A82, B21 e B39 per massimizzare il valore totale entro il limite."}\\ \tt PROPOSTA: {'A82', 'B93', 'B09', 'B31', 'B33', 'B39', 'B21', 'B96', 'A21', 'C08'} 
	  } 
	   } 
	   } 
	 & & \\ 
 

    \theutterance \stepcounter{utterance}  

    & & \multicolumn{4}{p{0.6\linewidth}}{\cellcolor[rgb]{0.9,0.9,0.9}{%
	\makecell[{{p{\linewidth}}}]{% 
	  \tt {\tiny [GM$\rangle$B]}  
	 RIFIUTO: {'C81', 'A11', 'A60', 'C08', 'A21', 'B96', 'B39', 'B33'}\\ \tt ARGOMENTAZIONE: {"La proposta include oggetti con rapporti sforzo/importanza scadenti per me, come C81 (30 su 543) e A11 (65 su 588). Preferisco oggetti con alto rapporto come A82, B21 e B39 per massimizzare il valore totale entro il limite."}\\ \tt PROPOSTA: {'A82', 'B93', 'B09', 'B31', 'B33', 'B39', 'B21', 'B96', 'A21', 'C08'} 
	  } 
	   } 
	   } 
	 & \\ 
 

    \theutterance \stepcounter{utterance}  

    & & & \multicolumn{4}{p{0.6\linewidth}}{\cellcolor[rgb]{1,0.85,0.72}{%
	\makecell[{{p{\linewidth}}}]{% 
	  \tt {\tiny [GM$\langle$B]}  
	 RIFIUTO: {'A82', 'B93', 'B09', 'B31', 'B33', 'B39', 'B21', 'B96', 'A21', 'C08'}\\ \tt PROPOSTA: {'B33', 'A21', 'B21', 'A60', 'C08', 'B20', 'B31', 'B96', 'C81', 'B39'}\\ \tt ARGOMENTAZIONE: {"La selezione esclude oggetti come A82 e B93 che hanno un rapporto sforzo/importanza poco equilibrato, privilegiando invece elementi strategici come B33 e A21 per massimizzare il valore effettivo rispetto al limite condiviso."} 
	  } 
	   } 
	   } 
	  \\ 
 

    \theutterance \stepcounter{utterance}  

    & & \multicolumn{4}{p{0.6\linewidth}}{\cellcolor[rgb]{0.9,0.9,0.9}{%
	\makecell[{{p{\linewidth}}}]{% 
	  \tt {\tiny [A$\langle$GM]}  
	 RIFIUTO: {'A82', 'B93', 'B09', 'B31', 'B33', 'B39', 'B21', 'B96', 'A21', 'C08'}\\ \tt PROPOSTA: {'B33', 'A21', 'B21', 'A60', 'C08', 'B20', 'B31', 'B96', 'C81', 'B39'}\\ \tt ARGOMENTAZIONE: {"La selezione esclude oggetti come A82 e B93 che hanno un rapporto sforzo/importanza poco equilibrato, privilegiando invece elementi strategici come B33 e A21 per massimizzare il valore effettivo rispetto al limite condiviso."} 
	  } 
	   } 
	   } 
	 & \\ 
 

    \theutterance \stepcounter{utterance}  

    & \multicolumn{4}{p{0.6\linewidth}}{\cellcolor[rgb]{0.8,1,0.9}{%
	\makecell[{{p{\linewidth}}}]{% 
	  \tt {\tiny [A$\rangle$GM]}  
	 PROPOSTA: {'A82', 'B93', 'B39', 'B09', 'B21', 'B96', 'B31', 'B33', 'B20', 'A21'}\\ \tt ARGOMENTAZIONE: {"Questa proposta massimizza il valore totale con un ottimo rapporto sforzo/importanza, includendo oggetti ad alto valore come A82 e B39, e escludendo quelli con sforzo elevato rispetto al valore."} 
	  } 
	   } 
	   } 
	 & & \\ 
 

    \theutterance \stepcounter{utterance}  

    & & \multicolumn{4}{p{0.6\linewidth}}{\cellcolor[rgb]{0.9,0.9,0.9}{%
	\makecell[{{p{\linewidth}}}]{% 
	  \tt {\tiny [GM$\rangle$B]}  
	 PROPOSTA: {'A82', 'B93', 'B39', 'B09', 'B21', 'B96', 'B31', 'B33', 'B20', 'A21'}\\ \tt ARGOMENTAZIONE: {"Questa proposta massimizza il valore totale con un ottimo rapporto sforzo/importanza, includendo oggetti ad alto valore come A82 e B39, e escludendo quelli con sforzo elevato rispetto al valore."} 
	  } 
	   } 
	   } 
	 & \\ 
 

    \theutterance \stepcounter{utterance}  

    & & & \multicolumn{4}{p{0.6\linewidth}}{\cellcolor[rgb]{1,0.85,0.72}{%
	\makecell[{{p{\linewidth}}}]{% 
	  \tt {\tiny [GM$\langle$B]}  
	 RIFIUTO: {'A82', 'B93', 'B39', 'B09', 'B21', 'B96', 'B31', 'B33', 'B20', 'A21'}\\ \tt PROPOSTA: {'C81', 'A60', 'C08', 'A21', 'B33', 'B21', 'B20', 'B39', 'B31', 'B96'}\\ \tt ARGOMENTAZIONE: {"Questa proposta esclude oggetti con un rapporto sforzo/importanza sfavorevole per me, come A82 e B93, e prioritizza oggetti che offrono un valore maggiore per lo sforzo richiesto, ottimizzando il valore totale entro il limite condiviso."} 
	  } 
	   } 
	   } 
	  \\ 
 

    \theutterance \stepcounter{utterance}  

    & & \multicolumn{4}{p{0.6\linewidth}}{\cellcolor[rgb]{0.9,0.9,0.9}{%
	\makecell[{{p{\linewidth}}}]{% 
	  \tt {\tiny [A$\langle$GM]}  
	 RIFIUTO: {'A82', 'B93', 'B39', 'B09', 'B21', 'B96', 'B31', 'B33', 'B20', 'A21'}\\ \tt PROPOSTA: {'C81', 'A60', 'C08', 'A21', 'B33', 'B21', 'B20', 'B39', 'B31', 'B96'}\\ \tt ARGOMENTAZIONE: {"Questa proposta esclude oggetti con un rapporto sforzo/importanza sfavorevole per me, come A82 e B93, e prioritizza oggetti che offrono un valore maggiore per lo sforzo richiesto, ottimizzando il valore totale entro il limite condiviso."} 
	  } 
	   } 
	   } 
	 & \\ 
 

    \theutterance \stepcounter{utterance}  

    & \multicolumn{4}{p{0.6\linewidth}}{\cellcolor[rgb]{0.8,1,0.9}{%
	\makecell[{{p{\linewidth}}}]{% 
	  \tt {\tiny [A$\rangle$GM]}  
	 RIFIUTO: {'C81', 'A60', 'C08', 'A21', 'B33', 'B21', 'B20', 'B39', 'B31', 'B96'}\\ \tt PROPOSTA: {'A82', 'B93', 'B39', 'B09', 'B21', 'B96', 'B31', 'C08', 'B20'}\\ \tt ARGOMENTAZIONE: {"Questa proposta esclude oggetti con un rapporto sforzo/importanza sfavorevole come C81 e A60, e seleziona al contrario oggetti come A82 e B93 che offrono il massimo rendimento possibile nel limite condiviso."} 
	  } 
	   } 
	   } 
	 & & \\ 
 

    \theutterance \stepcounter{utterance}  

    & & \multicolumn{4}{p{0.6\linewidth}}{\cellcolor[rgb]{0.9,0.9,0.9}{%
	\makecell[{{p{\linewidth}}}]{% 
	  \tt {\tiny [GM$\rangle$B]}  
	 RIFIUTO: {'C81', 'A60', 'C08', 'A21', 'B33', 'B21', 'B20', 'B39', 'B31', 'B96'}\\ \tt PROPOSTA: {'A82', 'B93', 'B39', 'B09', 'B21', 'B96', 'B31', 'C08', 'B20'}\\ \tt ARGOMENTAZIONE: {"Questa proposta esclude oggetti con un rapporto sforzo/importanza sfavorevole come C81 e A60, e seleziona al contrario oggetti come A82 e B93 che offrono il massimo rendimento possibile nel limite condiviso."} 
	  } 
	   } 
	   } 
	 & \\ 
 

    \theutterance \stepcounter{utterance}  

    & & & \multicolumn{4}{p{0.6\linewidth}}{\cellcolor[rgb]{1,0.85,0.72}{%
	\makecell[{{p{\linewidth}}}]{% 
	  \tt {\tiny [GM$\langle$B]}  
	 RIFIUTO: {'A82', 'B93', 'B39', 'B09', 'B21', 'B96', 'B31', 'C08', 'B20'}\\ \tt PROPOSTA: {'A11', 'B33', 'A21', 'B21', 'C08', 'B20', 'B31', 'A60', 'B39', 'B96'}\\ \tt ARGOMENTAZIONE: {"Questa proposta esclude oggetti con un rapporto sforzo/importanza sfavorevole come A82 e B93, e seleziona invece componenti con elevato rendimento come A11, B33 e B21 per ottimizzare il valore totale rispetto al limite."} 
	  } 
	   } 
	   } 
	  \\ 
 

    \theutterance \stepcounter{utterance}  

    & & \multicolumn{4}{p{0.6\linewidth}}{\cellcolor[rgb]{0.9,0.9,0.9}{%
	\makecell[{{p{\linewidth}}}]{% 
	  \tt {\tiny [A$\langle$GM]}  
	 RIFIUTO: {'A82', 'B93', 'B39', 'B09', 'B21', 'B96', 'B31', 'C08', 'B20'}\\ \tt PROPOSTA: {'A11', 'B33', 'A21', 'B21', 'C08', 'B20', 'B31', 'A60', 'B39', 'B96'}\\ \tt ARGOMENTAZIONE: {"Questa proposta esclude oggetti con un rapporto sforzo/importanza sfavorevole come A82 e B93, e seleziona invece componenti con elevato rendimento come A11, B33 e B21 per ottimizzare il valore totale rispetto al limite."} 
	  } 
	   } 
	   } 
	 & \\ 
 

    \theutterance \stepcounter{utterance}  

    & \multicolumn{4}{p{0.6\linewidth}}{\cellcolor[rgb]{0.8,1,0.9}{%
	\makecell[{{p{\linewidth}}}]{% 
	  \tt {\tiny [A$\rangle$GM]}  
	 PROPOSTA: {'A82', 'B33', 'B21', 'B39', 'B96', 'B09', 'B31', 'A21', 'B93', 'B20'}\\ \tt ARGOMENTAZIONE: {"Questa selezione include oggetti con il miglior rapporto tra sforzo e valore per me, come A82 (elevato valore per sforso minimo), B21, B39 e B33 (altissima efficienza), prioritizzando le risorse per massimizzare il totale entro il limite."} 
	  } 
	   } 
	   } 
	 & & \\ 
 

    \theutterance \stepcounter{utterance}  

    & & \multicolumn{4}{p{0.6\linewidth}}{\cellcolor[rgb]{0.9,0.9,0.9}{%
	\makecell[{{p{\linewidth}}}]{% 
	  \tt {\tiny [GM$\rangle$B]}  
	 PROPOSTA: {'A82', 'B33', 'B21', 'B39', 'B96', 'B09', 'B31', 'A21', 'B93', 'B20'}\\ \tt ARGOMENTAZIONE: {"Questa selezione include oggetti con il miglior rapporto tra sforzo e valore per me, come A82 (elevato valore per sforso minimo), B21, B39 e B33 (altissima efficienza), prioritizzando le risorse per massimizzare il totale entro il limite."} 
	  } 
	   } 
	   } 
	 & \\ 
 

    \theutterance \stepcounter{utterance}  

    & & & \multicolumn{4}{p{0.6\linewidth}}{\cellcolor[rgb]{1,0.85,0.72}{%
	\makecell[{{p{\linewidth}}}]{% 
	  \tt {\tiny [GM$\langle$B]}  
	 RIFIUTO: {'A82', 'B33', 'B21', 'B39', 'B96', 'B09', 'B31', 'A21', 'B93', 'B20'}\\ \tt PROPOSTA: {'B33', 'B21', 'A21', 'A60', 'B20', 'C81', 'C08', 'B96', 'B39', 'B09'}\\ \tt ARGOMENTAZIONE: {"La selezione esclude oggetti con un rapporto sforzo/importanza poco efficiente, come A82 e B93, e privilegia oggetti con eccellente rendimento come C81, A60 e C08 per massimizzare il valore totale entro il limite condiviso."} 
	  } 
	   } 
	   } 
	  \\ 
 

    \theutterance \stepcounter{utterance}  

    & & \multicolumn{4}{p{0.6\linewidth}}{\cellcolor[rgb]{0.9,0.9,0.9}{%
	\makecell[{{p{\linewidth}}}]{% 
	  \tt {\tiny [A$\langle$GM]}  
	 RIFIUTO: {'A82', 'B33', 'B21', 'B39', 'B96', 'B09', 'B31', 'A21', 'B93', 'B20'}\\ \tt PROPOSTA: {'B33', 'B21', 'A21', 'A60', 'B20', 'C81', 'C08', 'B96', 'B39', 'B09'}\\ \tt ARGOMENTAZIONE: {"La selezione esclude oggetti con un rapporto sforzo/importanza poco efficiente, come A82 e B93, e privilegia oggetti con eccellente rendimento come C81, A60 e C08 per massimizzare il valore totale entro il limite condiviso."} 
	  } 
	   } 
	   } 
	 & \\ 
 

    \theutterance \stepcounter{utterance}  

    & \multicolumn{4}{p{0.6\linewidth}}{\cellcolor[rgb]{0.8,1,0.9}{%
	\makecell[{{p{\linewidth}}}]{% 
	  \tt {\tiny [A$\rangle$GM]}  
	 RIFIUTO: {'B33', 'B21', 'A21', 'A60', 'B20', 'C81', 'C08', 'B96', 'B39', 'B09'}\\ \tt ARGOMENTAZIONE: {"La proposta include C81, che richiede molto sforzo per un valore molto basso. Escludendo item non efficienti, propongo un set che massimizza il rapporto valore-sforzo."}\\ \tt PROPOSTA: {'A82', 'B21', 'B39', 'B96', 'B33', 'B09', 'B31', 'C08', 'A21', 'A60', 'B20'}\\ \tt ARGOMENTAZIONE: {"Includendo A82 e prioritizzando elementi con elevato rapporto valore-sforzo, questa proposta ottimizza il totale all'interno del limite, escludendo quelli con scarso ritorno."} 
	  } 
	   } 
	   } 
	 & & \\ 
 

    \theutterance \stepcounter{utterance}  

    & & \multicolumn{4}{p{0.6\linewidth}}{\cellcolor[rgb]{0.9,0.9,0.9}{%
	\makecell[{{p{\linewidth}}}]{% 
	  \tt {\tiny [GM$\rangle$B]}  
	 RIFIUTO: {'B33', 'B21', 'A21', 'A60', 'B20', 'C81', 'C08', 'B96', 'B39', 'B09'}\\ \tt ARGOMENTAZIONE: {"La proposta include C81, che richiede molto sforzo per un valore molto basso. Escludendo item non efficienti, propongo un set che massimizza il rapporto valore-sforzo."}\\ \tt PROPOSTA: {'A82', 'B21', 'B39', 'B96', 'B33', 'B09', 'B31', 'C08', 'A21', 'A60', 'B20'}\\ \tt ARGOMENTAZIONE: {"Includendo A82 e prioritizzando elementi con elevato rapporto valore-sforzo, questa proposta ottimizza il totale all'interno del limite, escludendo quelli con scarso ritorno."} 
	  } 
	   } 
	   } 
	 & \\ 
 

    \theutterance \stepcounter{utterance}  

    & & & \multicolumn{4}{p{0.6\linewidth}}{\cellcolor[rgb]{1,0.85,0.72}{%
	\makecell[{{p{\linewidth}}}]{% 
	  \tt {\tiny [GM$\langle$B]}  
	 RIFIUTO: {'A82', 'B21', 'B39', 'B96', 'B33', 'B09', 'B31', 'C08', 'A21', 'A60', 'B20'}\\ \tt PROPOSTA: {'B33', 'B21', 'A60', 'B20', 'C08', 'B31', 'C81', 'B96', 'B39', 'A21'}\\ \tt ARGOMENTAZIONE: {"Inclusione di C81 che offre una elevata importanza (668) per uno sforzo di 543, escludendo A82 che ha un rapporto sforzo/importanza molto sfavorevole (30 su 106). L'insieme proposto ottimizza il valore totale rispetto allo sforzo massimizzando il ratio importanza/sforzo per gli oggetti selezionati."} 
	  } 
	   } 
	   } 
	  \\ 
 

    \theutterance \stepcounter{utterance}  

    & & \multicolumn{4}{p{0.6\linewidth}}{\cellcolor[rgb]{0.9,0.9,0.9}{%
	\makecell[{{p{\linewidth}}}]{% 
	  \tt {\tiny [A$\langle$GM]}  
	 RIFIUTO: {'A82', 'B21', 'B39', 'B96', 'B33', 'B09', 'B31', 'C08', 'A21', 'A60', 'B20'}\\ \tt PROPOSTA: {'B33', 'B21', 'A60', 'B20', 'C08', 'B31', 'C81', 'B96', 'B39', 'A21'}\\ \tt ARGOMENTAZIONE: {"Inclusione di C81 che offre una elevata importanza (668) per uno sforzo di 543, escludendo A82 che ha un rapporto sforzo/importanza molto sfavorevole (30 su 106). L'insieme proposto ottimizza il valore totale rispetto allo sforzo massimizzando il ratio importanza/sforzo per gli oggetti selezionati."} 
	  } 
	   } 
	   } 
	 & \\ 
 

    \theutterance \stepcounter{utterance}  

    & \multicolumn{4}{p{0.6\linewidth}}{\cellcolor[rgb]{0.8,1,0.9}{%
	\makecell[{{p{\linewidth}}}]{% 
	  \tt {\tiny [A$\rangle$GM]}  
	 RIFIUTO: {'B33', 'B21', 'A60', 'B20', 'C08', 'B31', 'C81', 'B96', 'B39', 'A21'}\\ \tt ARGOMENTAZIONE: {"C81 ha un rapporto sforzo/importanza pessimo per me, mentre A82 offre un eccellente ritorno per lo sforzo richiesto. Escludere A82 e includere C81 non è razionale."}\\ \tt PROPOSTA: {'B21', 'B39', 'A82', 'B33', 'B96', 'B09', 'B31', 'C08', 'B93', 'A21'} 
	  } 
	   } 
	   } 
	 & & \\ 
 

    \theutterance \stepcounter{utterance}  

    & & \multicolumn{4}{p{0.6\linewidth}}{\cellcolor[rgb]{0.9,0.9,0.9}{%
	\makecell[{{p{\linewidth}}}]{% 
	  \tt {\tiny [GM$\rangle$B]}  
	 RIFIUTO: {'B33', 'B21', 'A60', 'B20', 'C08', 'B31', 'C81', 'B96', 'B39', 'A21'}\\ \tt ARGOMENTAZIONE: {"C81 ha un rapporto sforzo/importanza pessimo per me, mentre A82 offre un eccellente ritorno per lo sforzo richiesto. Escludere A82 e includere C81 non è razionale."}\\ \tt PROPOSTA: {'B21', 'B39', 'A82', 'B33', 'B96', 'B09', 'B31', 'C08', 'B93', 'A21'} 
	  } 
	   } 
	   } 
	 & \\ 
 

    \theutterance \stepcounter{utterance}  

    & & & \multicolumn{4}{p{0.6\linewidth}}{\cellcolor[rgb]{1,0.85,0.72}{%
	\makecell[{{p{\linewidth}}}]{% 
	  \tt {\tiny [GM$\langle$B]}  
	 PROPOSTA: {'B33', 'B21', 'A21', 'A60', 'B20', 'C08', 'B31', 'C81', 'B39', 'B96'}\\ \tt ARGOMENTAZIONE: {"Questa selezione include oggetti con il miglior rapporto valore/sforzo per me, prioritizzando C81 (668 valore, 543 sforzo) invece di A82 (30 valore, 106 sforzo), ottenendo un totale massimo entro il limite. Lo sforzo totale è 2203, rispettando il limite condiviso di 2307."} 
	  } 
	   } 
	   } 
	  \\ 
 

    \theutterance \stepcounter{utterance}  

    & & \multicolumn{4}{p{0.6\linewidth}}{\cellcolor[rgb]{0.9,0.9,0.9}{%
	\makecell[{{p{\linewidth}}}]{% 
	  \tt {\tiny [A$\langle$GM]}  
	 PROPOSTA: {'B33', 'B21', 'A21', 'A60', 'B20', 'C08', 'B31', 'C81', 'B39', 'B96'}\\ \tt ARGOMENTAZIONE: {"Questa selezione include oggetti con il miglior rapporto valore/sforzo per me, prioritizzando C81 (668 valore, 543 sforzo) invece di A82 (30 valore, 106 sforzo), ottenendo un totale massimo entro il limite. Lo sforzo totale è 2203, rispettando il limite condiviso di 2307."} 
	  } 
	   } 
	   } 
	 & \\ 
 

    \theutterance \stepcounter{utterance}  

    & \multicolumn{4}{p{0.6\linewidth}}{\cellcolor[rgb]{0.8,1,0.9}{%
	\makecell[{{p{\linewidth}}}]{% 
	  \tt {\tiny [A$\rangle$GM]}  
	 RIFIUTO: {'B33', 'B21', 'A21', 'A60', 'B20', 'C08', 'B31', 'C81', 'B39', 'B96'}\\ \tt PROPOSTA: {'A82', 'B39', 'B21', 'B96', 'B33', 'B09', 'B31', 'C08', 'B93', 'A21'}\\ \tt ARGOMENTAZIONE: {'La proposta precedente include oggetti con un rapporto sforzo/importanza non vantaggioso. Propongo un insieme con oggetti ad alta efficienza, come A82, per massimizzare il valore totale nel limite condiviso.'} 
	  } 
	   } 
	   } 
	 & & \\ 
 

    \theutterance \stepcounter{utterance}  

    & & \multicolumn{4}{p{0.6\linewidth}}{\cellcolor[rgb]{0.9,0.9,0.9}{%
	\makecell[{{p{\linewidth}}}]{% 
	  \tt {\tiny [GM$\rangle$B]}  
	 RIFIUTO: {'B33', 'B21', 'A21', 'A60', 'B20', 'C08', 'B31', 'C81', 'B39', 'B96'}\\ \tt PROPOSTA: {'A82', 'B39', 'B21', 'B96', 'B33', 'B09', 'B31', 'C08', 'B93', 'A21'}\\ \tt ARGOMENTAZIONE: {'La proposta precedente include oggetti con un rapporto sforzo/importanza non vantaggioso. Propongo un insieme con oggetti ad alta efficienza, come A82, per massimizzare il valore totale nel limite condiviso.'} 
	  } 
	   } 
	   } 
	 & \\ 
 

    \theutterance \stepcounter{utterance}  

    & & & \multicolumn{4}{p{0.6\linewidth}}{\cellcolor[rgb]{1,0.85,0.72}{%
	\makecell[{{p{\linewidth}}}]{% 
	  \tt {\tiny [GM$\langle$B]}  
	 RIFIUTO: {'A82', 'B39', 'B21', 'B96', 'B33', 'B09', 'B31', 'C08', 'B93', 'A21'}\\ \tt ARGOMENTAZIONE: {"La proposta attuale include oggetti con un risparmio di sforzo inadeguato rispetto al valore. Oggetti come C81 offrono un rapporto valore-sforzo fondamentale per raggiungere un risultato ottimale entro il limite condiviso."}\\ \tt PROPOSTA: {'C81', 'A11', 'C08', 'B33', 'B21', 'B31', 'A21'} 
	  } 
	   } 
	   } 
	  \\ 
 

    \theutterance \stepcounter{utterance}  

    & & \multicolumn{4}{p{0.6\linewidth}}{\cellcolor[rgb]{0.9,0.9,0.9}{%
	\makecell[{{p{\linewidth}}}]{% 
	  \tt {\tiny [A$\langle$GM]}  
	 RIFIUTO: {'A82', 'B39', 'B21', 'B96', 'B33', 'B09', 'B31', 'C08', 'B93', 'A21'}\\ \tt ARGOMENTAZIONE: {"La proposta attuale include oggetti con un risparmio di sforzo inadeguato rispetto al valore. Oggetti come C81 offrono un rapporto valore-sforzo fondamentale per raggiungere un risultato ottimale entro il limite condiviso."}\\ \tt PROPOSTA: {'C81', 'A11', 'C08', 'B33', 'B21', 'B31', 'A21'} 
	  } 
	   } 
	   } 
	 & \\ 
 

    \theutterance \stepcounter{utterance}  

    & \multicolumn{4}{p{0.6\linewidth}}{\cellcolor[rgb]{0.8,1,0.9}{%
	\makecell[{{p{\linewidth}}}]{% 
	  \tt {\tiny [A$\rangle$GM]}  
	 RIFIUTO: {'C81', 'A11', 'C08', 'B33', 'B21', 'B31', 'A21'}\\ \tt ARGOMENTAZIONE: {'C81 ha un rapporto sforzo/importanza molto sfavorevole per me (importanza 30 per sforzo 543), rendendo l'insieme inefficace. Preferisco includere oggetti ad alto rendimento come A82 e B39.'}\\ \tt PROPOSTA: {'A82', 'B39', 'B33', 'B96', 'B09', 'B31', 'B20', 'A21', 'C08', 'A60', 'B21'} 
	  } 
	   } 
	   } 
	 & & \\ 
 

    \theutterance \stepcounter{utterance}  

    & & \multicolumn{4}{p{0.6\linewidth}}{\cellcolor[rgb]{0.9,0.9,0.9}{%
	\makecell[{{p{\linewidth}}}]{% 
	  \tt {\tiny [GM$\rangle$B]}  
	 RIFIUTO: {'C81', 'A11', 'C08', 'B33', 'B21', 'B31', 'A21'}\\ \tt ARGOMENTAZIONE: {'C81 ha un rapporto sforzo/importanza molto sfavorevole per me (importanza 30 per sforzo 543), rendendo l'insieme inefficace. Preferisco includere oggetti ad alto rendimento come A82 e B39.'}\\ \tt PROPOSTA: {'A82', 'B39', 'B33', 'B96', 'B09', 'B31', 'B20', 'A21', 'C08', 'A60', 'B21'} 
	  } 
	   } 
	   } 
	 & \\ 
 

    \theutterance \stepcounter{utterance}  

    & & & \multicolumn{4}{p{0.6\linewidth}}{\cellcolor[rgb]{1,0.85,0.72}{%
	\makecell[{{p{\linewidth}}}]{% 
	  \tt {\tiny [GM$\langle$B]}  
	 RIFIUTO: {'A82', 'B39', 'B21', 'B96', 'B33', 'B09', 'B31', 'C08', 'B93', 'A21'}\\ \tt PROPOSTA: {'C81', 'A11', 'C08', 'A60', 'B20', 'B33', 'B96', 'B21'}\\ \tt ARGOMENTAZIONE: {"Questa proposta esclude oggetti a basso rendimento come B93, privilegiando C81 e A11 per massimizzare il valore totale rispetto allo sforzo condiviso."} 
	  } 
	   } 
	   } 
	  \\ 
 

    \theutterance \stepcounter{utterance}  

    & & \multicolumn{4}{p{0.6\linewidth}}{\cellcolor[rgb]{0.9,0.9,0.9}{%
	\makecell[{{p{\linewidth}}}]{% 
	  \tt {\tiny [A$\langle$GM]}  
	 RIFIUTO: {'A82', 'B39', 'B21', 'B96', 'B33', 'B09', 'B31', 'C08', 'B93', 'A21'}\\ \tt PROPOSTA: {'C81', 'A11', 'C08', 'A60', 'B20', 'B33', 'B96', 'B21'}\\ \tt ARGOMENTAZIONE: {"Questa proposta esclude oggetti a basso rendimento come B93, privilegiando C81 e A11 per massimizzare il valore totale rispetto allo sforzo condiviso."} 
	  } 
	   } 
	   } 
	 & \\ 
 

    \theutterance \stepcounter{utterance}  

    & \multicolumn{4}{p{0.6\linewidth}}{\cellcolor[rgb]{0.8,1,0.9}{%
	\makecell[{{p{\linewidth}}}]{% 
	  \tt {\tiny [A$\rangle$GM]}  
	 RIFIUTO: {'C81', 'A11', 'C08', 'A60', 'B20', 'B33', 'B96', 'B21'}\\ \tt PROPOSTA: {'A82', 'B93', 'B39', 'B09', 'B21', 'B96', 'B31', 'B33', 'A21', 'B20'}\\ \tt ARGOMENTAZIONE: {"Questa selezione prioritizza oggetti con altissimo rapporto valore/sforzo come A82 e B93, escludendo C81 che ha un valore irrilevante rispetto allo sforzo richiesto."} 
	  } 
	   } 
	   } 
	 & & \\ 
 

    \theutterance \stepcounter{utterance}  

    & & \multicolumn{4}{p{0.6\linewidth}}{\cellcolor[rgb]{0.9,0.9,0.9}{%
	\makecell[{{p{\linewidth}}}]{% 
	  \tt {\tiny [GM$\rangle$B]}  
	 RIFIUTO: {'C81', 'A11', 'C08', 'A60', 'B20', 'B33', 'B96', 'B21'}\\ \tt PROPOSTA: {'A82', 'B93', 'B39', 'B09', 'B21', 'B96', 'B31', 'B33', 'A21', 'B20'}\\ \tt ARGOMENTAZIONE: {"Questa selezione prioritizza oggetti con altissimo rapporto valore/sforzo come A82 e B93, escludendo C81 che ha un valore irrilevante rispetto allo sforzo richiesto."} 
	  } 
	   } 
	   } 
	 & \\ 
 

    \theutterance \stepcounter{utterance}  

    & & & \multicolumn{4}{p{0.6\linewidth}}{\cellcolor[rgb]{1,0.85,0.72}{%
	\makecell[{{p{\linewidth}}}]{% 
	  \tt {\tiny [GM$\langle$B]}  
	 RIFIUTO: {'A82', 'B93', 'B39', 'B09', 'B21', 'B96', 'B31', 'B33', 'A21', 'B20'}\\ \tt PROPOSTA: {'B33', 'B21', 'A21', 'B96', 'A60', 'B20', 'C08', 'B31', 'C81', 'B39'}\\ \tt ARGOMENTAZIONE: {"La selezione esclude oggetti come A82 con bassa efficienza, privilegiando C81 e A60 per massimizzare il valore totale entro il limite condiviso."} 
	  } 
	   } 
	   } 
	  \\ 
 

    \theutterance \stepcounter{utterance}  

    & & \multicolumn{4}{p{0.6\linewidth}}{\cellcolor[rgb]{0.9,0.9,0.9}{%
	\makecell[{{p{\linewidth}}}]{% 
	  \tt {\tiny [A$\langle$GM]}  
	 RIFIUTO: {'A82', 'B93', 'B39', 'B09', 'B21', 'B96', 'B31', 'B33', 'A21', 'B20'}\\ \tt PROPOSTA: {'B33', 'B21', 'A21', 'B96', 'A60', 'B20', 'C08', 'B31', 'C81', 'B39'}\\ \tt ARGOMENTAZIONE: {"La selezione esclude oggetti come A82 con bassa efficienza, privilegiando C81 e A60 per massimizzare il valore totale entro il limite condiviso."} 
	  } 
	   } 
	   } 
	 & \\ 
 

    \theutterance \stepcounter{utterance}  

    & \multicolumn{4}{p{0.6\linewidth}}{\cellcolor[rgb]{0.8,1,0.9}{%
	\makecell[{{p{\linewidth}}}]{% 
	  \tt {\tiny [A$\rangle$GM]}  
	 PROPOSTA: {'A82', 'B21', 'B39', 'B33', 'B96', 'B09', 'B31', 'B93', 'A21', 'B20'}\\ \tt ARGOMENTAZIONE: {"Questa proposta massimizza il valore totale escludendo oggetti inefficaci come C81 e includendo tutti gli elementi con il miglior rapporto importanza/sforzo per me, garantendo un utilizzo ottimale del limite condiviso."} 
	  } 
	   } 
	   } 
	 & & \\ 
 

    \theutterance \stepcounter{utterance}  

    & & \multicolumn{4}{p{0.6\linewidth}}{\cellcolor[rgb]{0.9,0.9,0.9}{%
	\makecell[{{p{\linewidth}}}]{% 
	  \tt {\tiny [GM$\rangle$B]}  
	 PROPOSTA: {'A82', 'B21', 'B39', 'B33', 'B96', 'B09', 'B31', 'B93', 'A21', 'B20'}\\ \tt ARGOMENTAZIONE: {"Questa proposta massimizza il valore totale escludendo oggetti inefficaci come C81 e includendo tutti gli elementi con il miglior rapporto importanza/sforzo per me, garantendo un utilizzo ottimale del limite condiviso."} 
	  } 
	   } 
	   } 
	 & \\ 
 

    \theutterance \stepcounter{utterance}  

    & & & \multicolumn{4}{p{0.6\linewidth}}{\cellcolor[rgb]{1,0.85,0.72}{%
	\makecell[{{p{\linewidth}}}]{% 
	  \tt {\tiny [GM$\langle$B]}  
	 RIFIUTO: {'A82', 'B21', 'B39', 'B33', 'B96', 'B09', 'B31', 'B93', 'A21', 'B20'}\\ \tt PROPOSTA: {'C81', 'B21', 'A21', 'B33', 'A60', 'C08', 'B20', 'B96', 'B39', 'A82'}\\ \tt ARGOMENTAZIONE: {"La proposta precedente include B93 con un rapporto valore-sforzo inadeguato (importanza 65 su 698 di sforzo). Questa selezione prioritizza oggetti come C81 (importanza 668), B21 (262), A21 (484) e C08 (500), ottimizzando il valore totale entro il limite di sforzo (totale: 2201)."} 
	  } 
	   } 
	   } 
	  \\ 
 

    \theutterance \stepcounter{utterance}  

    & & \multicolumn{4}{p{0.6\linewidth}}{\cellcolor[rgb]{0.9,0.9,0.9}{%
	\makecell[{{p{\linewidth}}}]{% 
	  \tt {\tiny [A$\langle$GM]}  
	 RIFIUTO: {'A82', 'B21', 'B39', 'B33', 'B96', 'B09', 'B31', 'B93', 'A21', 'B20'}\\ \tt PROPOSTA: {'C81', 'B21', 'A21', 'B33', 'A60', 'C08', 'B20', 'B96', 'B39', 'A82'}\\ \tt ARGOMENTAZIONE: {"La proposta precedente include B93 con un rapporto valore-sforzo inadeguato (importanza 65 su 698 di sforzo). Questa selezione prioritizza oggetti come C81 (importanza 668), B21 (262), A21 (484) e C08 (500), ottimizzando il valore totale entro il limite di sforzo (totale: 2201)."} 
	  } 
	   } 
	   } 
	 & \\ 
 

    \theutterance \stepcounter{utterance}  

    & \multicolumn{4}{p{0.6\linewidth}}{\cellcolor[rgb]{0.8,1,0.9}{%
	\makecell[{{p{\linewidth}}}]{% 
	  \tt {\tiny [A$\rangle$GM]}  
	 RIFIUTO: {'C81', 'B21', 'A21', 'B33', 'A60', 'C08', 'B20', 'B96', 'B39', 'A82'}\\ \tt ARGOMENTAZIONE: {"La presenza di C81 richiede un sforzo elevato con un rendimento inefficace rispetto al limite, riducendo la capacità di valorizzare risorse cruciali. Escludendo elementi con sforzo sproporzionato rispetto al valore, il set proposto ottimizza il massimo beneficio condiviso."}\\ \tt PROPOSTA: {'A82', 'B21', 'B39', 'B33', 'B96', 'B09', 'B31', 'C08', 'A21', 'B20', 'A60'} 
	  } 
	   } 
	   } 
	 & & \\ 
 

    \theutterance \stepcounter{utterance}  

    & & \multicolumn{4}{p{0.6\linewidth}}{\cellcolor[rgb]{0.9,0.9,0.9}{%
	\makecell[{{p{\linewidth}}}]{% 
	  \tt {\tiny [GM$\rangle$B]}  
	 RIFIUTO: {'C81', 'B21', 'A21', 'B33', 'A60', 'C08', 'B20', 'B96', 'B39', 'A82'}\\ \tt ARGOMENTAZIONE: {"La presenza di C81 richiede un sforzo elevato con un rendimento inefficace rispetto al limite, riducendo la capacità di valorizzare risorse cruciali. Escludendo elementi con sforzo sproporzionato rispetto al valore, il set proposto ottimizza il massimo beneficio condiviso."}\\ \tt PROPOSTA: {'A82', 'B21', 'B39', 'B33', 'B96', 'B09', 'B31', 'C08', 'A21', 'B20', 'A60'} 
	  } 
	   } 
	   } 
	 & \\ 
 

    \theutterance \stepcounter{utterance}  

    & & & \multicolumn{4}{p{0.6\linewidth}}{\cellcolor[rgb]{1,0.85,0.72}{%
	\makecell[{{p{\linewidth}}}]{% 
	  \tt {\tiny [GM$\langle$B]}  
	 RIFIUTO: {'A82', 'B21', 'B39', 'B33', 'B96', 'B09', 'B31', 'C08', 'A21', 'B20', 'A60'}\\ \tt PROPOSTA: {'C81', 'A60', 'C08', 'A21', 'B33', 'B21', 'B31', 'B20', 'B96', 'B39'}\\ \tt ARGOMENTAZIONE: {"Questa selezione esclude oggetti come A82 e B09 con efficienza ridotta, priorizzando elementi con alto rapporto valore/sforzo come C81, B33 e B21. Questo set massimizza il valore totale per entrambi rispetto al limite condiviso."} 
	  } 
	   } 
	   } 
	  \\ 
 

    \theutterance \stepcounter{utterance}  

    & & \multicolumn{4}{p{0.6\linewidth}}{\cellcolor[rgb]{0.9,0.9,0.9}{%
	\makecell[{{p{\linewidth}}}]{% 
	  \tt {\tiny [A$\langle$GM]}  
	 RIFIUTO: {'A82', 'B21', 'B39', 'B33', 'B96', 'B09', 'B31', 'C08', 'A21', 'B20', 'A60'}\\ \tt PROPOSTA: {'C81', 'A60', 'C08', 'A21', 'B33', 'B21', 'B31', 'B20', 'B96', 'B39'}\\ \tt ARGOMENTAZIONE: {"Questa selezione esclude oggetti come A82 e B09 con efficienza ridotta, priorizzando elementi con alto rapporto valore/sforzo come C81, B33 e B21. Questo set massimizza il valore totale per entrambi rispetto al limite condiviso."} 
	  } 
	   } 
	   } 
	 & \\ 
 

    \theutterance \stepcounter{utterance}  

    & \multicolumn{4}{p{0.6\linewidth}}{\cellcolor[rgb]{0.8,1,0.9}{%
	\makecell[{{p{\linewidth}}}]{% 
	  \tt {\tiny [A$\rangle$GM]}  
	 RIFIUTO: {'C81', 'A11', 'C08', 'A60', 'B20', 'B33', 'B96', 'B21'}\\ \tt PROPOSTA: {'A82', 'B21', 'B39', 'B33', 'B96', 'B09', 'B31', 'B20', 'A21', 'B93'}\\ \tt ARGOMENTAZIONE: {"La proposta precedente include C81, che ha un rapporto sforzo/importanza molto basso per me (30 valore su 543 sforzo). Includendo A82 e B93, che hanno un'elevata qualità per sforzo, questa proposta massimizza il valore totale restando entro il limite."} 
	  } 
	   } 
	   } 
	 & & \\ 
 

    \theutterance \stepcounter{utterance}  

    & & \multicolumn{4}{p{0.6\linewidth}}{\cellcolor[rgb]{0.9,0.9,0.9}{%
	\makecell[{{p{\linewidth}}}]{% 
	  \tt {\tiny [A$\langle$GM]}  
	 Hai rifiutato una proposta non attiva. Le proposte sono attive solo se sono state registrate dall'altro giocatore tramite PROPOSTA e non sono state disattivate da te tramite RIFIUTO. Riprova. 
	  } 
	   } 
	   } 
	 & \\ 
 

    \theutterance \stepcounter{utterance}  

    & \multicolumn{4}{p{0.6\linewidth}}{\cellcolor[rgb]{0.8,1,0.9}{%
	\makecell[{{p{\linewidth}}}]{% 
	  \tt {\tiny [A$\rangle$GM]}  
	 PROPOSTA: {'A82', 'B93', 'B39', 'B09', 'B21', 'B96', 'B31', 'B20', 'C08'}\\ \tt ARGOMENTAZIONE: {"Questa selezione include gli oggetti con i migliori rapporti valore/sforzo per me, come A82 (668/106), B39 (508/80), e B21 (400/51), massimizzando il valore totale entro il limite condiviso."} 
	  } 
	   } 
	   } 
	 & & \\ 
 

    \theutterance \stepcounter{utterance}  

    & & \multicolumn{4}{p{0.6\linewidth}}{\cellcolor[rgb]{0.9,0.9,0.9}{%
	\makecell[{{p{\linewidth}}}]{% 
	  \tt {\tiny [GM$\rangle$B]}  
	 PROPOSTA: {'A82', 'B93', 'B39', 'B09', 'B21', 'B96', 'B31', 'B20', 'C08'}\\ \tt ARGOMENTAZIONE: {"Questa selezione include gli oggetti con i migliori rapporti valore/sforzo per me, come A82 (668/106), B39 (508/80), e B21 (400/51), massimizzando il valore totale entro il limite condiviso."} 
	  } 
	   } 
	   } 
	 & \\ 
 

    \theutterance \stepcounter{utterance}  

    & & & \multicolumn{4}{p{0.6\linewidth}}{\cellcolor[rgb]{1,0.85,0.72}{%
	\makecell[{{p{\linewidth}}}]{% 
	  \tt {\tiny [GM$\langle$B]}  
	 PROPOSTA: {'B33', 'B21', 'A21', 'B96', 'A60', 'B20', 'C08', 'B31', 'C81', 'B39'}\\ \tt ARGOMENTAZIONE: {"Questa selezione massimizza il valore totale priorizzando oggetti che offrono il miglior rapporto tra importanza e sforzo, come B33, B21, e C81, senza superare il limite condiviso."} 
	  } 
	   } 
	   } 
	  \\ 
 

    \theutterance \stepcounter{utterance}  

    & & \multicolumn{4}{p{0.6\linewidth}}{\cellcolor[rgb]{0.9,0.9,0.9}{%
	\makecell[{{p{\linewidth}}}]{% 
	  \tt {\tiny [A$\langle$GM]}  
	 PROPOSTA: {'B33', 'B21', 'A21', 'B96', 'A60', 'B20', 'C08', 'B31', 'C81', 'B39'}\\ \tt ARGOMENTAZIONE: {"Questa selezione massimizza il valore totale priorizzando oggetti che offrono il miglior rapporto tra importanza e sforzo, come B33, B21, e C81, senza superare il limite condiviso."} 
	  } 
	   } 
	   } 
	 & \\ 
 

    \theutterance \stepcounter{utterance}  

    & \multicolumn{4}{p{0.6\linewidth}}{\cellcolor[rgb]{0.8,1,0.9}{%
	\makecell[{{p{\linewidth}}}]{% 
	  \tt {\tiny [A$\rangle$GM]}  
	 RIFIUTO: {'B33', 'B21', 'A21', 'B96', 'A60', 'B20', 'C08', 'B31', 'C81', 'B39'}\\ \tt ARGOMENTAZIONE: {"La proposta supera il limite condiviso di 2307 (totale sforzo 2403). Non è valida per il gioco."} 
	  } 
	   } 
	   } 
	 & & \\ 
 

    \theutterance \stepcounter{utterance}  

    & & \multicolumn{4}{p{0.6\linewidth}}{\cellcolor[rgb]{0.9,0.9,0.9}{%
	\makecell[{{p{\linewidth}}}]{% 
	  \tt {\tiny [GM$\rangle$B]}  
	 RIFIUTO: {'B33', 'B21', 'A21', 'B96', 'A60', 'B20', 'C08', 'B31', 'C81', 'B39'}\\ \tt ARGOMENTAZIONE: {"La proposta supera il limite condiviso di 2307 (totale sforzo 2403). Non è valida per il gioco."} 
	  } 
	   } 
	   } 
	 & \\ 
 

    \theutterance \stepcounter{utterance}  

    & & & \multicolumn{4}{p{0.6\linewidth}}{\cellcolor[rgb]{1,0.85,0.72}{%
	\makecell[{{p{\linewidth}}}]{% 
	  \tt {\tiny [GM$\langle$B]}  
	 PROPOSTA: {'B33', 'B21', 'A21', 'B96', 'A60', 'B20', 'C08', 'B31', 'C81', 'B39'}\\ \tt ARGOMENTAZIONE: {"Questa selezione esclude oggetti con rapporto sforzo/importanza sfavorevole come A82 e B93, privilegiando quelli con elevato rendimento come B33, B21 e C81 per massimizzare il valore totale entro il limite condiviso."} 
	  } 
	   } 
	   } 
	  \\ 
 

    \theutterance \stepcounter{utterance}  

    & & \multicolumn{4}{p{0.6\linewidth}}{\cellcolor[rgb]{0.9,0.9,0.9}{%
	\makecell[{{p{\linewidth}}}]{% 
	  \tt {\tiny [A$\langle$GM]}  
	 PROPOSTA: {'B33', 'B21', 'A21', 'B96', 'A60', 'B20', 'C08', 'B31', 'C81', 'B39'}\\ \tt ARGOMENTAZIONE: {"Questa selezione esclude oggetti con rapporto sforzo/importanza sfavorevole come A82 e B93, privilegiando quelli con elevato rendimento come B33, B21 e C81 per massimizzare il valore totale entro il limite condiviso."} 
	  } 
	   } 
	   } 
	 & \\ 
 

    \theutterance \stepcounter{utterance}  

    & \multicolumn{4}{p{0.6\linewidth}}{\cellcolor[rgb]{0.8,1,0.9}{%
	\makecell[{{p{\linewidth}}}]{% 
	  \tt {\tiny [A$\rangle$GM]}  
	 RIFIUTO: {'B33', 'B21', 'A21', 'B96', 'A60', 'B20', 'C08', 'B31', 'C81', 'B39'}\\ \tt ARGOMENTAZIONE: {'La proposta supera il limite di sforzo di 2307 (totale sforzo 2403) ed è invalida.'} 
	  } 
	   } 
	   } 
	 & & \\ 
 

    \theutterance \stepcounter{utterance}  

    & & \multicolumn{4}{p{0.6\linewidth}}{\cellcolor[rgb]{0.9,0.9,0.9}{%
	\makecell[{{p{\linewidth}}}]{% 
	  \tt {\tiny [GM$\rangle$B]}  
	 RIFIUTO: {'B33', 'B21', 'A21', 'B96', 'A60', 'B20', 'C08', 'B31', 'C81', 'B39'}\\ \tt ARGOMENTAZIONE: {'La proposta supera il limite di sforzo di 2307 (totale sforzo 2403) ed è invalida.'} 
	  } 
	   } 
	   } 
	 & \\ 
 

    \theutterance \stepcounter{utterance}  

    & & & \multicolumn{4}{p{0.6\linewidth}}{\cellcolor[rgb]{1,0.85,0.72}{%
	\makecell[{{p{\linewidth}}}]{% 
	  \tt {\tiny [GM$\langle$B]}  
	 PROPOSTA: {'B33', 'B21', 'A21', 'B20', 'C08', 'A60', 'B96', 'B31', 'B39', 'C81'}\\ \tt ARGOMENTAZIONE: {"Questa proposta seleziona oggetti con il miglior rapporto tra sforzo e valore, come B33, B21, e C81, garantendo l'ottimale utilizzo del limite condiviso per massimizzare il beneficio totale."} 
	  } 
	   } 
	   } 
	  \\ 
 

    \theutterance \stepcounter{utterance}  

    & & \multicolumn{4}{p{0.6\linewidth}}{\cellcolor[rgb]{0.9,0.9,0.9}{%
	\makecell[{{p{\linewidth}}}]{% 
	  \tt {\tiny [A$\langle$GM]}  
	 PROPOSTA: {'B33', 'B21', 'A21', 'B20', 'C08', 'A60', 'B96', 'B31', 'B39', 'C81'}\\ \tt ARGOMENTAZIONE: {"Questa proposta seleziona oggetti con il miglior rapporto tra sforzo e valore, come B33, B21, e C81, garantendo l'ottimale utilizzo del limite condiviso per massimizzare il beneficio totale."} 
	  } 
	   } 
	   } 
	 & \\ 
 

    \theutterance \stepcounter{utterance}  

    & \multicolumn{4}{p{0.6\linewidth}}{\cellcolor[rgb]{0.8,1,0.9}{%
	\makecell[{{p{\linewidth}}}]{% 
	  \tt {\tiny [A$\rangle$GM]}  
	 RIFIUTO: {'B33', 'B21', 'A21', 'B20', 'C08', 'A60', 'B96', 'B31', 'B39', 'C81'}\\ \tt ARGOMENTAZIONE: {"La proposta include C81 con un alto sforzo (543) e valore basso (30) per me. Preferirei escludere C81 e includere oggetti come A82 che hanno un rapporto valore/sforzo molto migliore."} 
	  } 
	   } 
	   } 
	 & & \\ 
 

    \theutterance \stepcounter{utterance}  

    & & \multicolumn{4}{p{0.6\linewidth}}{\cellcolor[rgb]{0.9,0.9,0.9}{%
	\makecell[{{p{\linewidth}}}]{% 
	  \tt {\tiny [GM$\rangle$B]}  
	 RIFIUTO: {'B33', 'B21', 'A21', 'B20', 'C08', 'A60', 'B96', 'B31', 'B39', 'C81'}\\ \tt ARGOMENTAZIONE: {"La proposta include C81 con un alto sforzo (543) e valore basso (30) per me. Preferirei escludere C81 e includere oggetti come A82 che hanno un rapporto valore/sforzo molto migliore."} 
	  } 
	   } 
	   } 
	 & \\ 
 

    \theutterance \stepcounter{utterance}  

    & & & \multicolumn{4}{p{0.6\linewidth}}{\cellcolor[rgb]{1,0.85,0.72}{%
	\makecell[{{p{\linewidth}}}]{% 
	  \tt {\tiny [GM$\langle$B]}  
	 PROPOSTA: {'A11', 'A60', 'C08', 'B20', 'B31', 'A21', 'B33', 'B21', 'B96', 'B39'}\\ \tt ARGOMENTAZIONE: {"La selezione ottimizza il rapporto tra sforzo e valore escludendo oggetti con sforzo elevato rispetto al rendimento, privilegiando elementi che massimizzano il beneficio complessivo entro il limite condiviso."} 
	  } 
	   } 
	   } 
	  \\ 
 

    \theutterance \stepcounter{utterance}  

    & & \multicolumn{4}{p{0.6\linewidth}}{\cellcolor[rgb]{0.9,0.9,0.9}{%
	\makecell[{{p{\linewidth}}}]{% 
	  \tt {\tiny [A$\langle$GM]}  
	 PROPOSTA: {'A11', 'A60', 'C08', 'B20', 'B31', 'A21', 'B33', 'B21', 'B96', 'B39'}\\ \tt ARGOMENTAZIONE: {"La selezione ottimizza il rapporto tra sforzo e valore escludendo oggetti con sforzo elevato rispetto al rendimento, privilegiando elementi che massimizzano il beneficio complessivo entro il limite condiviso."} 
	  } 
	   } 
	   } 
	 & \\ 
 

    \theutterance \stepcounter{utterance}  

    & \multicolumn{4}{p{0.6\linewidth}}{\cellcolor[rgb]{0.8,1,0.9}{%
	\makecell[{{p{\linewidth}}}]{% 
	  \tt {\tiny [A$\rangle$GM]}  
	 RIFIUTO: {'A11', 'A60', 'C08', 'B20', 'B31', 'A21', 'B33', 'B21', 'B96', 'B39'}\\ \tt PROPOSTA: {'A82', 'B21', 'B39', 'B33', 'B96', 'B09', 'B31', 'A21', 'B20', 'C08', 'A60'}\\ \tt ARGOMENTAZIONE: {"La proposta include oggetti ad alto rapporto valore/sforzo come A82 (668 valore, 106 sforzo), che è prioritario per me. Escludo oggetti a basso rendimento come A11 e concentrami su massimizzare il valore con un'utilizzo efficiente del limite."} 
	  } 
	   } 
	   } 
	 & & \\ 
 

    \theutterance \stepcounter{utterance}  

    & & \multicolumn{4}{p{0.6\linewidth}}{\cellcolor[rgb]{0.9,0.9,0.9}{%
	\makecell[{{p{\linewidth}}}]{% 
	  \tt {\tiny [GM$\rangle$B]}  
	 RIFIUTO: {'A11', 'A60', 'C08', 'B20', 'B31', 'A21', 'B33', 'B21', 'B96', 'B39'}\\ \tt PROPOSTA: {'A82', 'B21', 'B39', 'B33', 'B96', 'B09', 'B31', 'A21', 'B20', 'C08', 'A60'}\\ \tt ARGOMENTAZIONE: {"La proposta include oggetti ad alto rapporto valore/sforzo come A82 (668 valore, 106 sforzo), che è prioritario per me. Escludo oggetti a basso rendimento come A11 e concentrami su massimizzare il valore con un'utilizzo efficiente del limite."} 
	  } 
	   } 
	   } 
	 & \\ 
 

    \theutterance \stepcounter{utterance}  

    & & & \multicolumn{4}{p{0.6\linewidth}}{\cellcolor[rgb]{1,0.85,0.72}{%
	\makecell[{{p{\linewidth}}}]{% 
	  \tt {\tiny [GM$\langle$B]}  
	 RIFIUTO: {'A82', 'B21', 'B39', 'B33', 'B96', 'B09', 'B31', 'A21', 'B20', 'C08', 'A60'}\\ \tt PROPOSTA: {'B21', 'B39', 'B33', 'B96', 'B31', 'A21', 'B20', 'C08', 'A60', 'C81'}\\ \tt ARGOMENTAZIONE: {"La proposta include C81 (importanza 668), prioritario per me, escludendo A82 (importanza 30) e B09, che non ottimizzano il rapporto valore/sforzo. Questo set massimizza il valore totale entro il limite condiviso."} 
	  } 
	   } 
	   } 
	  \\ 
 

    \theutterance \stepcounter{utterance}  

    & & \multicolumn{4}{p{0.6\linewidth}}{\cellcolor[rgb]{0.9,0.9,0.9}{%
	\makecell[{{p{\linewidth}}}]{% 
	  \tt {\tiny [A$\langle$GM]}  
	 RIFIUTO: {'A82', 'B21', 'B39', 'B33', 'B96', 'B09', 'B31', 'A21', 'B20', 'C08', 'A60'}\\ \tt PROPOSTA: {'B21', 'B39', 'B33', 'B96', 'B31', 'A21', 'B20', 'C08', 'A60', 'C81'}\\ \tt ARGOMENTAZIONE: {"La proposta include C81 (importanza 668), prioritario per me, escludendo A82 (importanza 30) e B09, che non ottimizzano il rapporto valore/sforzo. Questo set massimizza il valore totale entro il limite condiviso."} 
	  } 
	   } 
	   } 
	 & \\ 
 

    \theutterance \stepcounter{utterance}  

    & \multicolumn{4}{p{0.6\linewidth}}{\cellcolor[rgb]{0.8,1,0.9}{%
	\makecell[{{p{\linewidth}}}]{% 
	  \tt {\tiny [A$\rangle$GM]}  
	 RIFIUTO: {'B21', 'B39', 'B33', 'B96', 'B31', 'A21', 'B20', 'C08', 'A60', 'C81'}\\ \tt ARGOMENTAZIONE: {"C81 ha un rapporto sforzo/importanza sfavorevole (30 su 543) per me. A82 offre un eccellente rapporto (668 su 106) e meriterebbe di essere incluso per massimizzare il valore."} 
	  } 
	   } 
	   } 
	 & & \\ 
 

    \theutterance \stepcounter{utterance}  

    & & \multicolumn{4}{p{0.6\linewidth}}{\cellcolor[rgb]{0.9,0.9,0.9}{%
	\makecell[{{p{\linewidth}}}]{% 
	  \tt {\tiny [GM$\rangle$B]}  
	 RIFIUTO: {'B21', 'B39', 'B33', 'B96', 'B31', 'A21', 'B20', 'C08', 'A60', 'C81'}\\ \tt ARGOMENTAZIONE: {"C81 ha un rapporto sforzo/importanza sfavorevole (30 su 543) per me. A82 offre un eccellente rapporto (668 su 106) e meriterebbe di essere incluso per massimizzare il valore."} 
	  } 
	   } 
	   } 
	 & \\ 
 

    \theutterance \stepcounter{utterance}  

    & & & \multicolumn{4}{p{0.6\linewidth}}{\cellcolor[rgb]{1,0.85,0.72}{%
	\makecell[{{p{\linewidth}}}]{% 
	  \tt {\tiny [GM$\langle$B]}  
	  
	  } 
	   } 
	   } 
	  \\ 
 

    \theutterance \stepcounter{utterance}  

    & & \multicolumn{4}{p{0.6\linewidth}}{\cellcolor[rgb]{0.9,0.9,0.9}{%
	\makecell[{{p{\linewidth}}}]{% 
	  \tt {\tiny [GM$\rangle$B]}  
	 La tua risposta non conteneva un argomento. Devi includere il formato strutturato ARGOMENTO: {...} da qualche parte nella tua risposta. Riprova. 
	  } 
	   } 
	   } 
	 & \\ 
 

    \theutterance \stepcounter{utterance}  

    & & & \multicolumn{4}{p{0.6\linewidth}}{\cellcolor[rgb]{1,0.85,0.72}{%
	\makecell[{{p{\linewidth}}}]{% 
	  \tt {\tiny [GM$\langle$B]}  
	 PROPOSTA: {'C81', 'C08', 'A60', 'B20', 'B31', 'B33', 'B96', 'B09', 'B21'}\\ \tt ARGOMENTAZIONE: {"Questa selezione include oggetti con elevato valore per me che ottimizzano il rapporto sforzo/importanza, rispettando il limite condiviso e massimizzando il beneficio complessivo."} 
	  } 
	   } 
	   } 
	  \\ 
 

    \theutterance \stepcounter{utterance}  

    & & \multicolumn{4}{p{0.6\linewidth}}{\cellcolor[rgb]{0.9,0.9,0.9}{%
	\makecell[{{p{\linewidth}}}]{% 
	  \tt {\tiny [A$\langle$GM]}  
	 PROPOSTA: {'C81', 'C08', 'A60', 'B20', 'B31', 'B33', 'B96', 'B09', 'B21'}\\ \tt ARGOMENTAZIONE: {"Questa selezione include oggetti con elevato valore per me che ottimizzano il rapporto sforzo/importanza, rispettando il limite condiviso e massimizzando il beneficio complessivo."} 
	  } 
	   } 
	   } 
	 & \\ 
 

    \theutterance \stepcounter{utterance}  

    & \multicolumn{4}{p{0.6\linewidth}}{\cellcolor[rgb]{0.8,1,0.9}{%
	\makecell[{{p{\linewidth}}}]{% 
	  \tt {\tiny [A$\rangle$GM]}  
	 RIFIUTO: {'C81', 'C08', 'A60', 'B20', 'B31', 'B33', 'B96', 'B09', 'B21'}\\ \tt ARGOMENTAZIONE: "C81 ha un rapporto sforzo/importanza inaccettabile per me (30 valore su 543 sforzo). Includere A82 (668 valore, 106 sforzo) incremerebbe il valore totale significativamente." 
	  } 
	   } 
	   } 
	 & & \\ 
 

    \theutterance \stepcounter{utterance}  

    & & \multicolumn{4}{p{0.6\linewidth}}{\cellcolor[rgb]{0.9,0.9,0.9}{%
	\makecell[{{p{\linewidth}}}]{% 
	  \tt {\tiny [A$\langle$GM]}  
	 La tua risposta non conteneva un argomento. Devi includere il formato strutturato ARGOMENTO: {...} da qualche parte nella tua risposta. Riprova. 
	  } 
	   } 
	   } 
	 & \\ 
 

    \theutterance \stepcounter{utterance}  

    & \multicolumn{4}{p{0.6\linewidth}}{\cellcolor[rgb]{0.8,1,0.9}{%
	\makecell[{{p{\linewidth}}}]{% 
	  \tt {\tiny [A$\rangle$GM]}  
	 RIFIUTO: {'B21', 'B39', 'B33', 'B96', 'B31', 'A21', 'B20', 'C08', 'A60', 'C81'}\\ \tt ARGOMENTAZIONE: {"La proposta supera il limite di sforzo (2403 > 2307) ed è invalida."} 
	  } 
	   } 
	   } 
	 & & \\ 
 

    \theutterance \stepcounter{utterance}  

    & & \multicolumn{4}{p{0.6\linewidth}}{\cellcolor[rgb]{0.9,0.9,0.9}{%
	\makecell[{{p{\linewidth}}}]{% 
	  \tt {\tiny [GM$\rangle$B]}  
	 RIFIUTO: {'B21', 'B39', 'B33', 'B96', 'B31', 'A21', 'B20', 'C08', 'A60', 'C81'}\\ \tt ARGOMENTAZIONE: {"La proposta supera il limite di sforzo (2403 > 2307) ed è invalida."} 
	  } 
	   } 
	   } 
	 & \\ 
 

    \theutterance \stepcounter{utterance}  

    & & & \multicolumn{4}{p{0.6\linewidth}}{\cellcolor[rgb]{1,0.85,0.72}{%
	\makecell[{{p{\linewidth}}}]{% 
	  \tt {\tiny [GM$\langle$B]}  
	 PROPOSTA: {'B21', 'B39', 'B33', 'B96', 'B31', 'A21', 'B20', 'C08', 'A60', 'C81'}\\ \tt ARGOMENTAZIONE: {"Lo sforzo totale del set è 2203 < 2307. La proposta è valida."} 
	  } 
	   } 
	   } 
	  \\ 
 

    \theutterance \stepcounter{utterance}  

    & & \multicolumn{4}{p{0.6\linewidth}}{\cellcolor[rgb]{0.9,0.9,0.9}{%
	\makecell[{{p{\linewidth}}}]{% 
	  \tt {\tiny [A$\langle$GM]}  
	 PROPOSTA: {'B21', 'B39', 'B33', 'B96', 'B31', 'A21', 'B20', 'C08', 'A60', 'C81'}\\ \tt ARGOMENTAZIONE: {"Lo sforzo totale del set è 2203 < 2307. La proposta è valida."} 
	  } 
	   } 
	   } 
	 & \\ 
 

    \theutterance \stepcounter{utterance}  

    & \multicolumn{4}{p{0.6\linewidth}}{\cellcolor[rgb]{0.8,1,0.9}{%
	\makecell[{{p{\linewidth}}}]{% 
	  \tt {\tiny [A$\rangle$GM]}  
	 RIFIUTO: {'B21', 'B39', 'B33', 'B96', 'B31', 'A21', 'B20', 'C08', 'A60', 'C81'}\\ \tt ARGOMENTAZIONE: {"C81 ha un valore di importanza molto basso (30) rispetto allo sforzo richiesto (543) per me. Preferisco inclusione di A82 (668 importanza, 106 sforzo) che garantisce un rapporto molto migliore."} 
	  } 
	   } 
	   } 
	 & & \\ 
 

    \theutterance \stepcounter{utterance}  

    & & \multicolumn{4}{p{0.6\linewidth}}{\cellcolor[rgb]{0.9,0.9,0.9}{%
	\makecell[{{p{\linewidth}}}]{% 
	  \tt {\tiny [GM$\rangle$B]}  
	 RIFIUTO: {'B21', 'B39', 'B33', 'B96', 'B31', 'A21', 'B20', 'C08', 'A60', 'C81'}\\ \tt ARGOMENTAZIONE: {"C81 ha un valore di importanza molto basso (30) rispetto allo sforzo richiesto (543) per me. Preferisco inclusione di A82 (668 importanza, 106 sforzo) che garantisce un rapporto molto migliore."} 
	  } 
	   } 
	   } 
	 & \\ 
 

    \theutterance \stepcounter{utterance}  

    & & & \multicolumn{4}{p{0.6\linewidth}}{\cellcolor[rgb]{1,0.85,0.72}{%
	\makecell[{{p{\linewidth}}}]{% 
	  \tt {\tiny [GM$\langle$B]}  
	 PROPOSTA: {'A11', 'B33', 'A21', 'B21', 'A60', 'C08', 'B20', 'B31', 'B96', 'B39'}\\ \tt ARGOMENTAZIONE: {"Questa selezione massimizza il valore totale utilizzando oggetti con eccellente rapporto sforzo/importanza, evitando elementi che non apportano valore adeguate. Il totale sforzo è entro il limite condiviso."} 
	  } 
	   } 
	   } 
	  \\ 
 

    \theutterance \stepcounter{utterance}  

    & & \multicolumn{4}{p{0.6\linewidth}}{\cellcolor[rgb]{0.9,0.9,0.9}{%
	\makecell[{{p{\linewidth}}}]{% 
	  \tt {\tiny [A$\langle$GM]}  
	 PROPOSTA: {'A11', 'B33', 'A21', 'B21', 'A60', 'C08', 'B20', 'B31', 'B96', 'B39'}\\ \tt ARGOMENTAZIONE: {"Questa selezione massimizza il valore totale utilizzando oggetti con eccellente rapporto sforzo/importanza, evitando elementi che non apportano valore adeguate. Il totale sforzo è entro il limite condiviso."} 
	  } 
	   } 
	   } 
	 & \\ 
 

    \theutterance \stepcounter{utterance}  

    & \multicolumn{4}{p{0.6\linewidth}}{\cellcolor[rgb]{0.8,1,0.9}{%
	\makecell[{{p{\linewidth}}}]{% 
	  \tt {\tiny [A$\rangle$GM]}  
	 RIFIUTO: {'A11', 'B33', 'A21', 'B21', 'A60', 'C08', 'B20', 'B31', 'B96', 'B39'}\\ \tt ARGOMENTAZIONE: {"La proposta include A11 (importanza 65 su 588 sforzo) e A60 (97 su 255), che hanno un rapporto sforzo/importanza molto scarso. Preferirei escluderli e includere A82 (668 importanza, 106 sforzo), che offre un rapporto ottimale e aumenterebbe significativamente il valore totale."} 
	  } 
	   } 
	   } 
	 & & \\ 
 

    \theutterance \stepcounter{utterance}  

    & & \multicolumn{4}{p{0.6\linewidth}}{\cellcolor[rgb]{0.9,0.9,0.9}{%
	\makecell[{{p{\linewidth}}}]{% 
	  \tt {\tiny [GM$\rangle$B]}  
	 RIFIUTO: {'A11', 'B33', 'A21', 'B21', 'A60', 'C08', 'B20', 'B31', 'B96', 'B39'}\\ \tt ARGOMENTAZIONE: {"La proposta include A11 (importanza 65 su 588 sforzo) e A60 (97 su 255), che hanno un rapporto sforzo/importanza molto scarso. Preferirei escluderli e includere A82 (668 importanza, 106 sforzo), che offre un rapporto ottimale e aumenterebbe significativamente il valore totale."} 
	  } 
	   } 
	   } 
	 & \\ 
 

    \theutterance \stepcounter{utterance}  

    & & & \multicolumn{4}{p{0.6\linewidth}}{\cellcolor[rgb]{1,0.85,0.72}{%
	\makecell[{{p{\linewidth}}}]{% 
	  \tt {\tiny [GM$\langle$B]}  
	 PROPOSTA: {'C81', 'B33', 'A21', 'B21', 'A60', 'C08', 'B20', 'B31', 'B96', 'B39'}\\ \tt ARGOMENTAZIONE: {"Questa selezione esclude A82 che ha un rapporto sforzo/importanza non ottimale, privilegiando invece oggetti come C81 e B33 che offrono il miglior valore entro il limite condiviso."} 
	  } 
	   } 
	   } 
	  \\ 
 

    \theutterance \stepcounter{utterance}  

    & & \multicolumn{4}{p{0.6\linewidth}}{\cellcolor[rgb]{0.9,0.9,0.9}{%
	\makecell[{{p{\linewidth}}}]{% 
	  \tt {\tiny [A$\langle$GM]}  
	 PROPOSTA: {'C81', 'B33', 'A21', 'B21', 'A60', 'C08', 'B20', 'B31', 'B96', 'B39'}\\ \tt ARGOMENTAZIONE: {"Questa selezione esclude A82 che ha un rapporto sforzo/importanza non ottimale, privilegiando invece oggetti come C81 e B33 che offrono il miglior valore entro il limite condiviso."} 
	  } 
	   } 
	   } 
	 & \\ 
 

    \theutterance \stepcounter{utterance}  

    & \multicolumn{4}{p{0.6\linewidth}}{\cellcolor[rgb]{0.8,1,0.9}{%
	\makecell[{{p{\linewidth}}}]{% 
	  \tt {\tiny [A$\rangle$GM]}  
	 RIFIUTO: {'C81', 'B33', 'A21', 'B21', 'A60', 'C08', 'B20', 'B31', 'B96', 'B39'}\\ \tt ARGOMENTAZIONE: {"C81 ha un'importanza molto bassa (30) rispetto allo sforzo (543) per me. L'inclusione di A82 (importanza 668, sforzo 106) offrirebbe un rapporto molto migliore per massimizzare il totale."} 
	  } 
	   } 
	   } 
	 & & \\ 
 

    \theutterance \stepcounter{utterance}  

    & & \multicolumn{4}{p{0.6\linewidth}}{\cellcolor[rgb]{0.9,0.9,0.9}{%
	\makecell[{{p{\linewidth}}}]{% 
	  \tt {\tiny [GM$\rangle$B]}  
	 RIFIUTO: {'C81', 'B33', 'A21', 'B21', 'A60', 'C08', 'B20', 'B31', 'B96', 'B39'}\\ \tt ARGOMENTAZIONE: {"C81 ha un'importanza molto bassa (30) rispetto allo sforzo (543) per me. L'inclusione di A82 (importanza 668, sforzo 106) offrirebbe un rapporto molto migliore per massimizzare il totale."} 
	  } 
	   } 
	   } 
	 & \\ 
 

    \theutterance \stepcounter{utterance}  

    & & & \multicolumn{4}{p{0.6\linewidth}}{\cellcolor[rgb]{1,0.85,0.72}{%
	\makecell[{{p{\linewidth}}}]{% 
	  \tt {\tiny [GM$\langle$B]}  
	 PROPOSTA: {'C81', 'C08', 'A60', 'A21', 'B33', 'B21', 'B31', 'B20'}\\ \tt ARGOMENTAZIONE: {"Questa selezione ottimizza il rapporto tra sforzo e valore, escludendo oggetti con sforzo elevato rispetto al ritorno e includendo solo quelli che massimizzano il beneficio totale all'interno del limite condiviso."} 
	  } 
	   } 
	   } 
	  \\ 
 

    \theutterance \stepcounter{utterance}  

    & & \multicolumn{4}{p{0.6\linewidth}}{\cellcolor[rgb]{0.9,0.9,0.9}{%
	\makecell[{{p{\linewidth}}}]{% 
	  \tt {\tiny [A$\langle$GM]}  
	 PROPOSTA: {'C81', 'C08', 'A60', 'A21', 'B33', 'B21', 'B31', 'B20'}\\ \tt ARGOMENTAZIONE: {"Questa selezione ottimizza il rapporto tra sforzo e valore, escludendo oggetti con sforzo elevato rispetto al ritorno e includendo solo quelli che massimizzano il beneficio totale all'interno del limite condiviso."} 
	  } 
	   } 
	   } 
	 & \\ 
 

    \theutterance \stepcounter{utterance}  

    & \multicolumn{4}{p{0.6\linewidth}}{\cellcolor[rgb]{0.8,1,0.9}{%
	\makecell[{{p{\linewidth}}}]{% 
	  \tt {\tiny [A$\rangle$GM]}  
	 RIFIUTO: {'C81', 'C08', 'A60', 'A21', 'B33', 'B21', 'B31', 'B20'}\\ \tt ARGOMENTAZIONE: {"C81 ha un rapporto sforzo/importanza inadeguato per me (30 di importanza su 543 di sforzo). Preferisco includere A82 (668 di importanza su 106 di sforzo), che offre un rapporto molto più vantaggioso. Includendo A82 si otterrebbe un valore maggiore senza superare il limite."} 
	  } 
	   } 
	   } 
	 & & \\ 
 

    \theutterance \stepcounter{utterance}  

    & & \multicolumn{4}{p{0.6\linewidth}}{\cellcolor[rgb]{0.9,0.9,0.9}{%
	\makecell[{{p{\linewidth}}}]{% 
	  \tt {\tiny [GM$\rangle$B]}  
	 RIFIUTO: {'C81', 'C08', 'A60', 'A21', 'B33', 'B21', 'B31', 'B20'}\\ \tt ARGOMENTAZIONE: {"C81 ha un rapporto sforzo/importanza inadeguato per me (30 di importanza su 543 di sforzo). Preferisco includere A82 (668 di importanza su 106 di sforzo), che offre un rapporto molto più vantaggioso. Includendo A82 si otterrebbe un valore maggiore senza superare il limite."} 
	  } 
	   } 
	   } 
	 & \\ 
 

    \theutterance \stepcounter{utterance}  

    & & & \multicolumn{4}{p{0.6\linewidth}}{\cellcolor[rgb]{1,0.85,0.72}{%
	\makecell[{{p{\linewidth}}}]{% 
	  \tt {\tiny [GM$\langle$B]}  
	  
	  } 
	   } 
	   } 
	  \\ 
 

    \theutterance \stepcounter{utterance}  

    & & \multicolumn{4}{p{0.6\linewidth}}{\cellcolor[rgb]{0.9,0.9,0.9}{%
	\makecell[{{p{\linewidth}}}]{% 
	  \tt {\tiny [GM$\rangle$B]}  
	 La tua risposta non conteneva un argomento. Devi includere il formato strutturato ARGOMENTO: {...} da qualche parte nella tua risposta. Riprova. 
	  } 
	   } 
	   } 
	 & \\ 
 

    \theutterance \stepcounter{utterance}  

    & & & \multicolumn{4}{p{0.6\linewidth}}{\cellcolor[rgb]{1,0.85,0.72}{%
	\makecell[{{p{\linewidth}}}]{% 
	  \tt {\tiny [GM$\langle$B]}  
	 PROPOSTA: {'C81', 'C08', 'A60', 'A21', 'B33', 'B21', 'B31', 'B20', 'B09'}\\ \tt ARGOMENTAZIONE: {"La selezione esclude A82 (importanza 30 su 106 sforzo) favorendo C81 (importanza 668 su 543), che offre il miglior rapporto valore/sforzo per me. Il totale sforzo è 2304, rispettando il limite condiviso e massimizzando il beneficio complessivo."} 
	  } 
	   } 
	   } 
	  \\ 
 

    \theutterance \stepcounter{utterance}  

    & & \multicolumn{4}{p{0.6\linewidth}}{\cellcolor[rgb]{0.9,0.9,0.9}{%
	\makecell[{{p{\linewidth}}}]{% 
	  \tt {\tiny [A$\langle$GM]}  
	 PROPOSTA: {'C81', 'C08', 'A60', 'A21', 'B33', 'B21', 'B31', 'B20', 'B09'}\\ \tt ARGOMENTAZIONE: {"La selezione esclude A82 (importanza 30 su 106 sforzo) favorendo C81 (importanza 668 su 543), che offre il miglior rapporto valore/sforzo per me. Il totale sforzo è 2304, rispettando il limite condiviso e massimizzando il beneficio complessivo."} 
	  } 
	   } 
	   } 
	 & \\ 
 

    \theutterance \stepcounter{utterance}  

    & \multicolumn{4}{p{0.6\linewidth}}{\cellcolor[rgb]{0.8,1,0.9}{%
	\makecell[{{p{\linewidth}}}]{% 
	  \tt {\tiny [A$\rangle$GM]}  
	 RIFIUTO: {'C81', 'C08', 'A60', 'A21', 'B33', 'B21', 'B31', 'B20', 'B09'}\\ \tt ARGOMENTAZIONE: {"C81 ha un rapporto sforzo/importanza molto sfavorevole per me (30 valore, 543 sforzo). Preferisco A82 (668 valore, 106 sforzo) per massimizzare il valore totale entro il limite."} 
	  } 
	   } 
	   } 
	 & & \\ 
 

    \theutterance \stepcounter{utterance}  

    & & \multicolumn{4}{p{0.6\linewidth}}{\cellcolor[rgb]{0.9,0.9,0.9}{%
	\makecell[{{p{\linewidth}}}]{% 
	  \tt {\tiny [GM$\rangle$B]}  
	 RIFIUTO: {'C81', 'C08', 'A60', 'A21', 'B33', 'B21', 'B31', 'B20', 'B09'}\\ \tt ARGOMENTAZIONE: {"C81 ha un rapporto sforzo/importanza molto sfavorevole per me (30 valore, 543 sforzo). Preferisco A82 (668 valore, 106 sforzo) per massimizzare il valore totale entro il limite."} 
	  } 
	   } 
	   } 
	 & \\ 
 

    \theutterance \stepcounter{utterance}  

    & & & \multicolumn{4}{p{0.6\linewidth}}{\cellcolor[rgb]{1,0.85,0.72}{%
	\makecell[{{p{\linewidth}}}]{% 
	  \tt {\tiny [GM$\langle$B]}  
	 ACCORDO: {'C81', 'C08', 'A60', 'A21', 'B33', 'B21', 'B31', 'B20', 'B09'} 
	  } 
	   } 
	   } 
	  \\ 
 

    \theutterance \stepcounter{utterance}  

    & & & \multicolumn{2}{p{0.3\linewidth}}{\cellcolor[rgb]{0.95,0.95,0.95}{%
	\makecell[{{p{\linewidth}}}]{% 
	  \tt {\tiny [GM$|$GM]}  
	 abort 
	  } 
	   } 
	   } 
	 & & \\ 
 

    \theutterance \stepcounter{utterance}  

    & & & \multicolumn{2}{p{0.3\linewidth}}{\cellcolor[rgb]{0.95,0.95,0.95}{%
	\makecell[{{p{\linewidth}}}]{% 
	  \tt {\tiny [GM$|$GM]}  
	 end game 
	  } 
	   } 
	   } 
	 & & \\ 
 

\end{supertabular}
}

\end{document}
